% concatenated from main14-arxiv1.tex
\documentclass[11pt,a4paper]{amsart}

\usepackage{amsmath,amssymb,amsthm,mathtools}
\usepackage[hidelinks]{hyperref}
\usepackage{mathrsfs}

\newtheorem{theorem}{Theorem}[section]
\newtheorem{proposition}[theorem]{Proposition}
\newtheorem{lemma}[theorem]{Lemma}
\newtheorem{corollary}[theorem]{Corollary}
\newtheorem*{maintheorem}{Theorem}
\theoremstyle{remark}
\newtheorem{remark}[theorem]{Remark}

\newcommand{\N}{\mathbb N}
\newcommand{\Z}{\mathbb Z}

\title[Sharpness of Bak--Sneppen avalanches]
{Sharpness of the avalanche phase transition in the Bak--Sneppen model}

% \title{Avalanches Thresholds in One-Dimensional Bak--Sneppen Models}

\keywords{Bak--Sneppen model, avalanche ranges, nonlinear jump processes, critical thresholds, self-organized criticality}

\author{Olivier Hénard}
\address{Institut de mathématique d'Orsay, Université Paris-Saclay,
91405 Orsay, France}
\email{olivier.henard@universite-paris-saclay.fr}

\hypersetup{
  pdftitle={Sharpness of the avalanche phase transition in the Bak--Sneppen model},
  pdfauthor={Olivier Hénard}
}

\begin{document}


\begin{abstract}
We prove sharpness of the avalanche phase transition in the
one-dimensional Bak--Sneppen model, for both the one-sided and isotropic
update rules. In the one-sided model, the avalanche range is described by
a nonlinear self-interacting jump process whose tails satisfy a triangular
system. An extremal inequality for decreasing probability masses yields a
polynomial lower bound on the range tail at the susceptibility threshold;
this bound forces the appearance of infinite avalanches immediately above
criticality and hence identifies the susceptibility and infinite-avalanche
thresholds. A comparison of avalanche endpoints transfers this equality to
the isotropic model. In both models, no infinite avalanche occurs at
criticality, while the range has exponential moments at every strictly
subcritical level.  Combined with the range--duration and
stationary-limit results of Meester and Znamenski, this identifies all three
avalanche critical thresholds and proves the conjectured stationary product
limit in both models.
\end{abstract}


\maketitle

\section{Introduction}
\label{sec:introduction}

The Bak--Sneppen model was introduced in \cite{BS93} as a simple model of
self-organized criticality driven by extremal dynamics. On the cycle
$\Lambda_N:=\mathbb Z/N\mathbb Z$, each site carries a fitness, and at
each step the least fit site together with a prescribed set of its nearest
neighbours is refreshed by independent uniform random variables. 
We consider the one-sided rule, also called the anisotropic rule in the physics
literature, which refreshes the minimum site and its right neighbour, and the
isotropic rule, which refreshes the minimum site and both neighbours.


We work throughout in exponential coordinates, obtained from the uniform
fitnesses by the order-preserving transformation
$u\mapsto-\log(1-u)$. Thus refreshed fitnesses are independent
$\operatorname{Exp}(1)$ variables. Throughout, $t\geq0$ denotes a fitness
level in these coordinates, not the discrete time of the Bak--Sneppen
dynamics.
For $\gamma\in\{\mathrm{os},\mathrm{iso}\}$, where $\mathrm{os}$ and
$\mathrm{iso}$ stand for the one-sided and isotropic update rules,
respectively, let $\pi_N^\gamma$ denote the unique stationary law of the
corresponding finite chain, and let
% For
% $\gamma\in\{\mathrm{os},\mathrm{iso}\}$, where $\mathrm{os}$ and
% $\mathrm{iso}$ stand for the one-sided and isotropic update rules,
% respectively, let $\pi_N^\gamma$ denote the
% stationary law of the corresponding chain, and let
\[
        F_N^\gamma
        =\bigl(F_{N,x}^\gamma\bigr)_{x\in\Lambda_N}
\]
have law $\pi_N^\gamma$.


The conjectured infinite-volume stationary picture is the emergence of a
deterministic fitness cutoff together with asymptotic independence for any
fixed collection of sites \cite{BS93,PMB96,MZ04}.
For $t\geq0$, let $\nu_t$ denote the exponential distribution conditioned
to lie above $t$, that is the law of $t+E$, where
$E\sim\operatorname{Exp}(1)$.

Following Meester and Znamenski \cite{MZ03,MZ04}, we study the stationary
problem through avalanches. The infinite-volume $t$-avalanche for update
rule $\gamma$ is obtained on $\mathbb Z$ by starting from a configuration
whose fitnesses all exceed $t$, forcing one update at the origin and its
neighbours prescribed by $\gamma$, and then running the usual dynamics
until all fitnesses again exceed $t$.
Let $\eta_\infty^\gamma(t)$ and
$r_\infty^\gamma(t)$ denote its duration and the number of sites in its
range, respectively, with both quantities set to $\infty$ if the avalanche
does not terminate. Define
\[
\begin{aligned}
        t_{\mathrm d}^\gamma
        &:=
        \sup\{t\geq0:
        \mathbb E[\eta_\infty^\gamma(t)]<\infty\},\\
        t_{\mathrm{sus}}^\gamma
        &:=
        \sup\{t\geq0:
        \mathbb E[r_\infty^\gamma(t)]<\infty\},\\
        t_\infty^\gamma
        &:=
        \inf\{t\geq0:
        \mathbb P(r_\infty^\gamma(t)=\infty)>0\}.
\end{aligned}
\]

The following statement collects the main conclusions of
Theorems~\ref{thm:range-criticality},
\ref{thm:isotropic-criticality}, and
\ref{thm:stationary-product-limits}.
% , together with
% Proposition~\ref{prop:isotropic-critical-finiteness}.


\begin{maintheorem}[Stationary product limit and avalanche sharpness]
For each $\gamma\in\{\mathrm{os},\mathrm{iso}\}$,
\[
        t_{\mathrm d}^\gamma
        =t_{\mathrm{sus}}^\gamma
        =t_\infty^\gamma
        =:t_c^\gamma.
\]
% Moreover,
% \[
%         \mathbb P\bigl(r^\gamma_\infty(t^\gamma_c)=\infty\bigr)=0,
%         \qquad
%         \mathbb P\bigl(r^\gamma_\infty(t^\gamma_c+s)=\infty\bigr)>0,
%         \qquad s>0,
% \]
% and, for every \(t<t^\gamma_c\), there exists \(\lambda_t>0\) such that \(
%         \mathbb E\!\left[e^{\lambda_t r^\gamma_\infty(t)}\right]<\infty.
% \) 
The critical values satisfy
\[
        1\leq t_c^{\mathrm{os}}\leq2\log2,
        \qquad
        \frac13\leq t_c^{\mathrm{iso}}
        \leq t_c^{\mathrm{os}}.
\]
Moreover, for every fixed collection of distinct sites
$x_1,\ldots,x_\ell\in\mathbb Z$, viewed as elements of $\Lambda_N$ for sufficiently large $N$,
\[
        \bigl(
        F_{N,x_1}^\gamma,\ldots,F_{N,x_\ell}^\gamma
        \bigr)
        \xrightarrow[N\to\infty]{\mathrm d}
        \nu_{t_c^\gamma}^{\otimes\ell}.
\]
Equivalently, in the original uniform parametrisation, the limiting
fitnesses are independent and uniform on $[q_c^\gamma,1]$, where
$q_c^\gamma:=1-e^{-t_c^\gamma}$.
\end{maintheorem}

The theorem closes, in the isotropic model, the programme initiated by
Meester and Znamenski. They first proved non-triviality of the large-system
stationary behaviour and developed a self-similar graphical representation
of avalanches \cite{MZ03}. They subsequently introduced the three
infinite-volume critical levels $t_{\mathrm d}^\gamma$, $t_{\mathrm{sus}}^\gamma$, 
        $t_\infty^\gamma$, proved equality of the first two and showed
that equality of all three implies the stationary product limit \cite{MZ04}. The
remaining problem was therefore a sharpness statement: divergence of the
mean range had to coincide with the first appearance of an infinite
avalanche.

Several related approaches and models were investigated while this
sharpness problem remained open. 
% Maslov first derived an infinite hierarchy of exact equations for avalanche
% observables \cite{Maslov96}.
Maslov, De Los Rios, Marsili and Zhang 
obtained an exact equation for the range-size distribution in the anisotropic
one-dimensional model, from which they extracted scaling relations and
critical exponents \cite{MDRMZ98}.
Rigorous comparison
arguments later gave upper and lower bounds on avalanche critical values on
transitive graphs by means of branching-process and percolation bounds
\cite{GMN06}, while so-called
maximal avalanches were shown to have
infinite expected durations in \cite{GMW06}.
Threshold limits have also been established for more tractable rank-driven
models and related topology-free surrogates of Bak--Sneppen dynamics
\cite{GKW11,GKW12}. For the Bernoulli-valued discrete Bak--Sneppen model,
a rigorous upper bound on the transition parameter was obtained by
renewal-theoretic methods \cite{Volkov22}. 

The proof separates the tractable one-sided update rule from the genuinely
two-sided isotropic update rule. In the one-sided model the avalanche range
closes in a scalar nonlinear evolution. Its tails satisfy a triangular
system closely related to the spatial-size equation of \cite{MDRMZ98}. An extremal inequality for decreasing probability masses controls the
second tail moment by the first and yields
\[
        \mathbb P(X_{t_c^{\mathrm{os}}}\geq k)
        \geq ck^{-2/3}.
\]
Only the fact that the exponent is strictly smaller than one is needed:
when the critical graft law is frozen, the expected times required to cross
successive dyadic scales are summable, so the frozen process explodes with positive probability in
arbitrarily small positive time. 
A triangular comparison then gives
one-sided sharpness.
% A stochastic ordering argument with the original process then gives
% one-sided sharpness.

The isotropic range has no closed scalar evolution. A single endpoint is,
however, an integral supersolution of the one-sided tail system. Restarting
the one-sided tail flow from a truncated endpoint law transfers the
preceding explosion mechanism and proves isotropic sharpness without
requiring a closed equation for the total range. Generating-function estimates give
exponential moments for the range at every strictly subcritical level for both
update rules. The final section and Appendix~\ref{app:mz-reduction} implement
the Meester--Znamenski range--duration and locking-threshold reduction in a
common form. 

\subsection*{Acknowledgements}
\leavevmode\par\noindent
The author thanks Linglong Yuan and Nicolas Fournier for helpful early
discussions.

\subsection*{Disclosure on the use of AI}
\leavevmode\par\noindent
Generative AI tools, including OpenAI's ChatGPT (GPT-5.6 Sol), were used during the exploratory stages of this work and in the preparation of the manuscript. In particular, AI-assisted exploration suggested the
extremal inequality stated in Lemma~\ref{lem:decreasing-mass}. 



\section{The one-sided range process}
\label{sec:avalanche-range}

We begin with the one-sided model, studied in \cite{MDRMZ98}, and use the exponential parametrisation. A fitness configuration is an element of $[0,\infty)^{\Z_+}$; at each discrete time, the site of minimum fitness and its right neighbour are refreshed by independent exponential random variables of mean one. A threshold $t$ in this parametrisation corresponds to $1-e^{-t}$ in the usual uniform parametrisation.

Fix $t\geq0$. Starting from a configuration whose fitnesses all exceed $t$, force one update of the sites $0$ and $1$, and thereafter update the site of minimum fitness and its right neighbour for as long as some fitness is at most $t$. The resulting procedure is the $t$-avalanche with origin $0$. Its duration $\eta(t)$ is the number of updates, including the forced update, and its range set $\xi(t)$ is the set of sites refreshed at least once. We write $r(t):=|\xi(t)|$, with the convention $r(t)=\infty$ when the range is infinite, and set $\eta(t)=\infty$ if the procedure never terminates. The law of $(\xi(t),\eta(t))$ does not depend on the initial values above $t$.


\begin{lemma}
\label{lem:range-interval}
For every $t\geq0$, almost surely,
\[
        \xi(t)=\{0,\ldots,r(t)-1\}
\]
when $r(t)<\infty$, whereas $\xi(t)=\mathbb Z_+$ when $r(t)=\infty$.
Moreover,
\[
        \{\eta(t)=\infty\}=\{r(t)=\infty\}
        \qquad\text{almost surely}.
\]
\end{lemma}

\begin{proof}
After the forced update the range is $\{0,1\}$. If the current range is
$\{0,\ldots,\ell\}$ and the selected site is $i$, the refreshed pair is
contained in $\{0,\ldots,\ell+1\}$. Thus the range remains an interval
with left endpoint $0$.

It remains to exclude infinite duration with finite range. On the event
that the range stays in $\{0,\ldots,L-1\}$, there are at most $L$ sites
below the threshold. If, during the next $L$ updates, both refreshed
fitnesses exceed $t$ at every update, then the number of sites below $t$
decreases at each step and the avalanche stops. This event has conditional
probability at least $e^{-2tL}$. Applying the estimate to successive
blocks of $L$ updates gives
\[
        \mathbb P\bigl(\eta(t)=\infty,\ r(t)\leq L\bigr)
        \leq\lim_{n\to\infty}(1-e^{-2tL})^n=0.
\]
Taking the union over $L$ proves the result.
\end{proof}




The avalanche family can be coupled monotonically in the threshold by the
self-similar graphical representation of Meester and Znamenski
\cite[Section~4]{MZ04}, with the updated neighbourhood
$\{-1,0,1\}$ replaced by $\{0,1\}$. We recall only the part needed below.
During the exploration of a $t$-avalanche, every refreshed value at most
$t$ is revealed exactly, whereas a value larger than $t$ is revealed only
through the event that it exceeds $t$. Let $\mathcal H_t$ be the
$\sigma$-field generated by the explored genealogy, the identities of the
refresh variables encountered, and this partially revealed information.
The family $(\mathcal H_t)_{t\geq0}$ is increasing.

\begin{lemma}[Regeneration at a level]
\label{lem:level-regeneration}
Suppose that the $t$-avalanche has finite range $I$. Conditionally on
$\mathcal H_t$, its terminal fitnesses are independent and have the form
\[
        F_x=t+E_x,\qquad x\in I,
\]
where the $E_x$ are independent $\operatorname{Exp}(1)$ variables,
independent of $\mathcal H_t$. Consequently, when the threshold is raised
above $t$, the next triggering level occurs at rate $|I|$, its location is
uniform on $I$, and the subavalanche attached there is independent of
$\mathcal H_t$ and has the law of an avalanche at the triggering level.
\end{lemma}

\begin{proof}
On $\{I\text{ is finite}\}$, the exploration determines, for every
$x\in I$, the identity of the last refresh variable assigned to $x$.
These variables are distinct and, conditionally on $\mathcal H_t$, each
has been observed only through the event that it exceeds $t$. Independence
and the lack-of-memory property therefore give the stated independent
excesses. Their minimum is exponential with rate $|I|$, and its location
is uniform on $I$. 
\end{proof}



Iterating the regeneration property gives the following self-similar
graphical representation. For $s\geq0$, set
$\mu_s:=\mathcal L(\xi(s))$. Independently for $x\in\Z_+$, take a marked
Poisson point process
\[
        \Pi_x=\sum_{j\geq1}\delta_{(\tau_{x,j},A_{x,j})}
        \quad\text{on }(0,\infty)\times2^{\Z_+}
        \quad\text{with intensity }ds\,\mu_s(dA).
\]
Thus $(\tau_{x,j})_{j\geq1}$ is a rate-one Poisson process and,
conditionally on $\tau_{x,j}=s$, the mark $A_{x,j}$ is an independent
copy of the range set of an $s$-avalanche rooted at the origin. At such
a mark, draw horizontal arrows from $(x,s)$ to $(x+y,s)$ for every
$y\in A_{x,j}$, and allow upward motion along the vertical lines.
This defines a random directed graph on $\Z_+ \times [0,\infty)$.
Starting from the seed $\{0,1\}$ at level $0$, let
$\xi^{\mathrm{gr}}(t)$ be the set of sites reachable at level $t$ by
directed paths of vertical segments and horizontal arrows.

For $n\geq3$, let
\[
        \sigma_n^{\mathrm{gr}}
        :=\inf\{t\geq0:|\xi^{\mathrm{gr}}(t)|\geq n\},
        \qquad
        \sigma_n:=\inf\{t\geq0:|\xi(t)|\geq n\}.
\]
Before either stopping level, fewer than $n$ vertical lines are active;
hence only finitely many reachable marks occur on every bounded level
interval. Lemma~\ref{lem:level-regeneration} may therefore be iterated
without accumulation and gives
\[
        \bigl(\xi^{\mathrm{gr}}
        (t\wedge\sigma_n^{\mathrm{gr}})\bigr)_{t\geq0}
        \stackrel{\mathrm d}=
        \bigl(\xi(t\wedge\sigma_n)\bigr)_{t\geq0}.
\]
Letting $n\to\infty$ and using monotonicity yields
\[
        (\xi^{\mathrm{gr}}(t))_{t\geq0}
        \stackrel{\mathrm d}=
        (\xi(t))_{t\geq0}.
\]
Since reachability is increasing in the terminal level, this
representation provides a simultaneous monotone coupling of the
avalanche ranges.

Define the shifted range process
\[
        X_t:=r(t)-1\in\N\cup\{\infty\},\qquad t\geq0.
\]
Thus $X_0=1$ and, on $\{X_t=x<\infty\}$, the current range is $\{0,\ldots,x\}$.




\begin{proposition}[Infinitesimal range evolution]
\label{prop:range-evolution}
At level $t$, conditionally on $X_t=x<\infty$, the update rate is
$x+1$, with update kernel
\begin{equation}
\label{eq:range-update}
        x\longmapsto\max\{x,U_x+X_t'\},
\end{equation}
where $U_x$ is uniform on $\{0,\ldots,x\}$ and $X_t'$ is independent of $U_x$ with the law of $X_t$. Null changes are allowed. The process is increasing and is sent to $\infty$ when its graphical range becomes infinite.
\end{proposition}

\begin{proof}
The superposition of the $x+1$ active rate-one Poisson processes gives the next mark at rate $x+1$, at a uniform location $U_x$. The mark carries an independent avalanche at the current level, whose shifted range is $\{0,\ldots,X_t'\}$. Its union with the current range is
\[
        \{0,\ldots,\max\{x,U_x+X_t'\}\},
\]
which proves the update rule.
\end{proof}

For $k\geq1$, set
\[
        q_k(t):=\mathbb{P}(X_t\geq k),\qquad
        q_\infty(t):=\mathbb{P}(X_t=\infty)=\lim_{k\to\infty}q_k(t).
\]
The next proposition is the only input from the avalanche construction used in the analytic part of the paper.

\begin{proposition}[Triangular tail system]
\label{prop:tail-system}
Each $q_k$ is non-decreasing and continuously differentiable in $t$,
while $k\mapsto q_k(t)$ is non-increasing. Moreover, $q_1\equiv1$,
$q_k(0)=0$ for $k\geq2$, and
\begin{equation}
\label{eq:tail-system}
        q_k'
        =
        \sum_{u=1}^{k-1}q_u(q_{k-u}-q_k)+q_k(1-q_k),
        \qquad k\geq2.
\end{equation}
In particular, the marginal law of the avalanche range is uniquely determined by this triangular system.
\end{proposition}

\begin{proof}
Fix $k\geq2$. Below $k$ the proposed rate is bounded by $k$, so the usual compensator calculation is legitimate. If $X_t=x<k$, then the rate at which a proposed update crosses level $k$ is
\[
        (x+1)\mathbb{P}(U_x+X_t'\geq k)
        =\sum_{j=0}^x q_{k-j}(t).
\]

Writing $p_x=q_x-q_{x+1}$, the compensator identity gives, for
$0\leq s\leq t$,
\[
\begin{split}
        q_k(t)-q_k(s)
        &=\int_s^t\sum_{x=1}^{k-1}p_x(u)
          \sum_{j=0}^xq_{k-j}(u)\,du\\
        &=\int_s^t\left[
          \sum_{v=1}^{k-1}q_v(u)(q_{k-v}(u)-q_k(u))
          +q_k(u)(1-q_k(u))\right]\,du.
\end{split}
\]
The initial conditions follow from $X_0=1$. Inductively in $k$, the integral
identity first shows that $q_k$ is continuous; its integrand is then
continuous, so $q_k$ is continuously differentiable and satisfies
\eqref{eq:tail-system} pointwise. At level $k$, this is a scalar locally
Lipschitz equation once $q_1,\ldots,q_{k-1}$ are known, and uniqueness of
the marginals follows by recursion.
\end{proof}

Write
\[
        m(t):=\mathbb{E}[X_t]=\sum_{k\geq1}q_k(t),\qquad
        t_{\mathrm{sus}}^{\mathrm{os}}
        :=\sup\{t\geq0:m(t)<\infty\},
\]
and
\[
        t_{\infty}^{\mathrm{os}}
        :=\inf\{t\geq0:\mathbb{P}(X_t=\infty)>0\}
        =\inf\{t\geq0:q_\infty(t)>0\}.
\]
Positive mass at infinity forces an infinite mean, so
$t_{\mathrm{sus}}^{\mathrm{os}}\leq t_{\infty}^{\mathrm{os}}$.

\begin{theorem}[One-sided avalanche range criticality]
\label{thm:range-criticality}
For the one-sided avalanche range process,
\[
        t_{\mathrm{sus}}^{\mathrm{os}}
        =t_{\infty}^{\mathrm{os}}
        =:t_c^{\mathrm{os}},
        \qquad 1\leq t_c^{\mathrm{os}}\leq2\log2.
\]
Moreover,
\[
        \mathbb{P}(X_{t_c^{\mathrm{os}}}=\infty)=0,
        \qquad
        \mathbb{P}(X_{t_c^{\mathrm{os}}+s}=\infty)>0,
        \quad s>0.
\]
There exists \(c>0\) such that
\[
        \mathbb{P}(X_{t_c^{\mathrm{os}}}\geq k)
        \geq ck^{-2/3},
        \qquad k\geq1.
\]
For every \(t<t_c^{\mathrm{os}}\), there exists \(\lambda_t>0\) such that
\(\mathbb{E}[e^{\lambda_tX_t}]<\infty\).
\end{theorem}

Set
\[
        t_*:=t_{\mathrm{sus}}^{\mathrm{os}}.
\]

\section{Tail evolution and susceptibility blow-up}
\label{sec:tail-evolution}

Equation~\eqref{eq:tail-system} may be written
\begin{equation}
\label{eq:tail-gain-loss}
        q_k'
        =
        \sum_{u=1}^{k-1}q_uq_{k-u}
        -q_k\sum_{u=2}^{k}q_u.
\end{equation}
Every term in \eqref{eq:tail-system} is non-negative because
\(k\mapsto q_k(t)\) is non-increasing. Consequently,
\begin{equation}
\label{eq:q-prime-gain-bound}
        0\leq q_k'(t)
        \leq
        g_k(t):=\sum_{u=1}^{k-1}q_u(t)q_{k-u}(t).
\end{equation}

\begin{lemma}[Differentiation of tail sums]
\label{lem:tail-differentiation}
Fix \(t_0<t_{\mathrm{sus}}^{\mathrm{os}}\). Then
\(\sum_{k\geq1}q_k'(t)\) converges uniformly on \([0,t_0]\), and
\[
        m'(t)=\sum_{k\geq1}q_k'(t).
\]
If \(A(t_0):=\sum_{k\geq1}kq_k(t_0)<\infty\), then
\(\sum_{k\geq1}kq_k'(t)\) also converges uniformly on \([0,t_0]\), and
\[
        A'(t)=\sum_{k\geq1}kq_k'(t).
\]
\end{lemma}

\begin{proof}
For \(0\leq t\leq t_0\), monotonicity in time and
\eqref{eq:q-prime-gain-bound} give
\(0\leq q_k'(t)\leq g_k(t_0)\), while
\[
        \sum_{k\geq2}g_k(t_0)=m(t_0)^2<\infty.
\]
The Weierstrass test proves the first assertion. If \(A(t_0)<\infty\),
then
\[
        \sum_{k\geq2}kg_k(t_0)
        =
        \sum_{u,v\geq1}(u+v)q_u(t_0)q_v(t_0)
        =
        2m(t_0)A(t_0)<\infty,
\]
and the second assertion follows in the same way.
\end{proof}

\begin{proposition}[First moment]
\label{prop:first-moment}
For \(t<t_*\),
\begin{equation}
\label{eq:mean-identity}
        m'
        =
        \frac12m^2+m-\frac12\sum_{k\geq1}q_k^2.
\end{equation}
In particular,
\begin{equation}
\label{eq:riccati-lower}
        m'\geq\frac12m(m+1),
\end{equation}
\[
        1\leq t_*\leq2\log2,
        \qquad
        \int_0^t(m(s)+1)\,ds\leq2\log m(t),\quad t<t_*,
\]
and \(m(t)\uparrow\infty\) as \(t\uparrow t_*\).
\end{proposition}

\begin{proof}
Let \(m_n:=\sum_{k=1}^nq_k\). From
\eqref{eq:q-prime-gain-bound}, \(m_n'\leq m_n^2\). Since \(m_n(0)=1\),
comparison with \(y'=y^2\) gives \(m_n(t)\leq(1-t)^{-1}\) for \(t<1\);
letting \(n\to\infty\) yields \(t_*\geq1\).

For \(t<t_*\), Lemma~\ref{lem:tail-differentiation} permits summation of
\eqref{eq:tail-gain-loss}. The gain is \(m^2\), whereas
\[
        \sum_{k\geq2}q_k\sum_{u=2}^{k}q_u
        =
        \frac12\left((m-1)^2+\sum_{k\geq2}q_k^2\right),
\]
which proves \eqref{eq:mean-identity}. Since
\(\sum_kq_k^2\leq m\), \eqref{eq:riccati-lower} follows from \eqref{eq:mean-identity}. Comparison with
\(y'=y(y+1)/2\), \(y(0)=1\), gives \(t_*\leq2\log2\), and
\((\log m)'\geq(m+1)/2\) gives the integral estimate.

Suppose finally that \(m(t)\uparrow M<\infty\) as \(t\uparrow t_*\).
Monotone convergence gives \(m(t_*)=M\). Integrating
\(m_n'\leq m_n^2\) from \(t_*\) yields
\[
        m_n(t_*+s)
        \leq
        \frac{m_n(t_*)}{1-sm_n(t_*)}
        \leq
        \frac{M}{1-sM},
        \qquad 0\leq s<M^{-1}.
\]
Letting \(n\to\infty\) contradicts the definition of \(t_*\).
\end{proof}

\begin{proposition}[Second moment below criticality]
\label{prop:polynomial-moments}
For every \(t<t_*\),
\[
        A(t):=\sum_{k\geq1}kq_k(t)<\infty,
        \qquad
        \mathbb{E}[X_t^2]=2A(t)-m(t)<\infty.
\]
\end{proposition}

\begin{proof}
Set \(A_n(t):=\sum_{k=1}^n kq_k(t)\). By
\eqref{eq:q-prime-gain-bound},
\[
        A_n'(t)
        \leq
        \sum_{\substack{u,v\geq1\\u+v\leq n}}
        (u+v)q_u(t)q_v(t)
        \leq2m(t)A_n(t).
\]
Hence
\[
        A_n(t)
        \leq
        \exp\left\{2\int_0^tm(s)\,ds\right\},
\]
uniformly in \(n\). The right-hand side is finite by
Proposition~\ref{prop:first-moment}; monotone convergence gives
\(A(t)<\infty\). The identity for the second moment follows from
\(X_t^2=\sum_{k=1}^{X_t}(2k-1)\) and Tonelli's theorem.
\end{proof}

\section{The extremal inequality and the critical tail}
\label{sec:critical-tail}

The central estimate is an extremal property of decreasing probability
masses.

\begin{lemma}[Mixtures of discrete uniform laws]
\label{lem:uniform-mixture}
Let \((p_k)_{k\geq1}\) be a non-increasing probability mass function and
set \(a_n:=n(p_n-p_{n+1})\). Then \((a_n)_{n\geq1}\) is a probability mass
function and
\begin{equation}
\label{eq:uniform-mixture}
        p_k=\sum_{n\geq k}\frac{a_n}{n}.
\end{equation}
Equivalently, if \(\mathbb{P}(N=n)=a_n\), then a random variable which,
conditionally on \(N=n\), is uniform on \(\{1,\ldots,n\}\), has law
\((p_k)\).
\end{lemma}

\begin{proof}
Monotonicity and summability imply \(np_n\to0\), since
\[
        np_n\leq2\sum_{j=\lceil n/2\rceil}^{n}p_j\longrightarrow0.
\]
Hence
\[
        \sum_{n=1}^Ma_n
        =
        \sum_{n=1}^Mn(p_n-p_{n+1})
        =
        \sum_{n=1}^Mp_n-Mp_{M+1}\longrightarrow1,
\]
and \(\sum_{n\geq k}a_n/n=\sum_{n\geq k}(p_n-p_{n+1})=p_k\).
\end{proof}

\begin{lemma}[Maximum under a decreasing mass function]
\label{lem:decreasing-mass}
Let \(K,K'\) be independent and identically distributed, with a non-increasing probability mass
function on \(\N\) and finite mean. Then
\begin{equation}
\label{eq:decreasing-max}
        \mathbb{E}[\max\{K,K'\}]
        \geq
        \frac43\mathbb{E}[K]-\frac13.
\end{equation}
\end{lemma}

\begin{proof}
Let \(U_a,U_b\) be independent and uniform on \(\{1,\ldots,a\}\) and
\(\{1,\ldots,b\}\), and assume \(a\leq b\). 
Then
\[
\begin{aligned}
ab\,\mathbb{E}[\max\{U_a,U_b\}]
&=\sum_{i=1}^a\left(i^2+\sum_{j=i+1}^b j\right) \\
&=\frac{b(b+1)}2\,a
  +\frac12\sum_{i=1}^a(i^2-i)
 =\frac{ab(b+1)}2+\frac{a(a^2-1)}6.
\end{aligned}
\]
Dividing by \(ab\) and rearranging gives
\[
        \mathbb{E}[\max\{U_a,U_b\}]
        =
        \frac{a+b+1}{3}
        +
        \frac{(b-a)^2+b-1}{6b} \geq\frac{a+b+1}{3} .
\]


Use Lemma~\ref{lem:uniform-mixture}, and let \(N,N'\) be independent with
law \((a_n)\). Conditionally on \((N,N')\), represent \(K,K'\) by
independent uniforms on \(\{1,\ldots,N\}\) and
\(\{1,\ldots,N'\}\). Then
\[
        \mathbb{E}[\max\{K,K'\}]
        \geq\frac{2\mathbb{E}[N]+1}{3}
        =
        \frac43\mathbb{E}[K]-\frac13,
\]
because \(\mathbb{E}[K]=(\mathbb{E}[N]+1)/2\).
\end{proof}

\begin{proposition}[Compensation estimate]
\label{prop:moment-comparison}
For every \(t<t_*\),
\begin{equation}
\label{eq:moment-comparison}
        A(t)\leq m(t)^{8/3}.
\end{equation}
\end{proposition}

\begin{proof}
By Proposition~\ref{prop:polynomial-moments} and
Lemma~\ref{lem:tail-differentiation}, the weighted tail equation may be
summed term by term. Put
\[
        L:=\sum_{1\leq u<k}kq_uq_k,
        \qquad
        D:=\sum_{k\geq1}kq_k^2.
\]
Summing \eqref{eq:tail-system} gives
\begin{equation}
\label{eq:A-prime}
        A'=2mA-L+A-D.
\end{equation}
Let \(K,K'\) be independent with law
\((q_k/m)_{k \geq 1}\). Then
\[
        \mathbb{E}[K]=\frac{A}{m},
        \qquad
        m^2\mathbb{E}[\max\{K,K'\}]=2L+D.
\]
Lemma~\ref{lem:decreasing-mass} therefore yields
\begin{equation}
\label{eq:L-compensation}
        L
        \geq
        \frac23mA-\frac16m^2-\frac12D.
\end{equation}
Combining \eqref{eq:A-prime} and \eqref{eq:L-compensation},
\[
        A'
        \leq
        \frac43mA+\frac16m^2+A-\frac12D.
\]
Since \(m^2\leq\mathbb{E}[X_t^2]=2A-m\),
\[
        A'
        \leq
        \frac43mA+\frac43A-\frac16m-\frac12D
        \leq
        \frac43(m+1)A.
\]
By \eqref{eq:riccati-lower},
\[
        \frac{A'}A
        \leq\frac43(m+1)
        \leq\frac83\frac{m'}m.
\]
Since \(A(0)=m(0)=1\), integration proves
\eqref{eq:moment-comparison}.
\end{proof}

\begin{proposition}[Critical lower tail]
\label{prop:critical-lower-tail}
There exists \(c>0\) such that
\[
        q_k(t_*)=\mathbb{P}(X_{t_*}\geq k)\geq ck^{-2/3},
        \qquad k\geq1.
\]
In particular, \(\mathbb{E}[X_{t_*}] = \sum_{k\geq1}q_k(t_*)=\infty\).
\end{proposition}

\begin{proof}
By Proposition~\ref{prop:moment-comparison},
\(\mathbb{E}[X_t^2]=2A(t)-m(t)\leq2m(t)^{8/3}\). Paley--Zygmund gives
\[
        \mathbb{P}\!\left(X_t\geq\frac12m(t)\right)
        \geq
        \frac14\frac{m(t)^2}{\mathbb{E}[X_t^2]}
        \geq\frac18m(t)^{-2/3}.
\]
By Proposition~\ref{prop:first-moment}, \(m\) is continuous on
\([0,t_*)\) and tends to infinity at \(t_*\). For every sufficiently large
integer \(n\), choose \(t_n<t_*\) such that \(m(t_n)=2n\). Monotonicity in
time gives
\[
        q_n(t_*)\geq q_n(t_n)\geq\frac18(2n)^{-2/3}.
\]
The triangular system implies inductively that \(q_n(t)>0\) for every
\(n\geq2\) and \(t>0\); decreasing the constant proves the estimate for
all \(n\geq1\).
\end{proof}

\begin{proposition}[No infinite avalanche at criticality]
\label{prop:critical-no-infinity}
At the susceptibility threshold,
\[
        q_\infty(t_*)=\mathbb{P}(X_{t_*}=\infty)=0.
\]
Thus $X_{t_*}<\infty$ almost surely.
\end{proposition}

Let us point this continuity statement is not logically needed to derive the 
equality of one-sided critical thresholds in Proposition~\ref{prop:equality-critical-times}.






\begin{proof}
Set $a:=q_\infty(t_*)$. Since by \eqref{eq:riccati-lower} $m'\geq m^2/2$ on $[0,t_*)$ and
$m(t)\uparrow\infty$ as $t\uparrow t_*$, integration of
$(1/m)'\leq-1/2$ gives
\begin{equation}
\label{eq:critical-mean-upper}
        m(t)\leq\frac{2}{t_*-t},
        \qquad t<t_*.
\end{equation}
Fix $C>0$, and for all sufficiently large $n$ set
$t_n:=t_*-C/n$. Since \(k\mapsto q_k(t_n)\) is non-increasing, 
$nq_n(t_n)\leq m(t_n)\leq2n/C$, and hence
\begin{equation}
\label{eq:critical-diagonal-upper}
        q_n(t_n)\leq\frac{2}{C}.
\end{equation}

We claim that $q_n(t_n)\to a$. For $s\in[t_n,t_*]$, one has
$q_{n-u}(s)\geq q_n(s)$, while monotonicity in the level gives
$q_{n-u}(s)\leq q_{n-u}(t_*)$. Hence
\eqref{eq:tail-system}, $q_u(s)\leq1$, and
$q_n(s)(1-q_n(s))\leq1$ give
\[
\begin{aligned}
        q_n'(s)
        &\leq
        \sum_{u=1}^{n-1}\bigl(q_{n-u}(t_*)-q_n(s)\bigr)+1\\
        &=B_n-(n-1)q_n(s),
        \qquad
        B_n:=1+\sum_{j=1}^{n-1}q_j(t_*).
\end{aligned}
\]
Writing $b_n:=B_n/(n-1)$ and integrating this scalar differential
inequality from $t_n$ to $t_*$ yields
\[
        q_n(t_*)
        \leq
        e^{-(n-1)C/n}q_n(t_n)
        +b_n\bigl(1-e^{-(n-1)C/n}\bigr).
\]
Now $q_n(t_*)\to a$, while $b_n\to a$ by Ces\`aro convergence.
Rearranging the preceding inequality gives
$\liminf_n q_n(t_n)\geq a$. Conversely,
$q_n(t_n)\leq q_n(t_*)$, so $\limsup_n q_n(t_n)\leq a$.
Thus $q_n(t_n)\to a$.

Combining this convergence with \eqref{eq:critical-diagonal-upper}
gives $a\leq2/C$. Since $C>0$ is arbitrary, $a=0$. 
\end{proof}

\section{Frozen critical dynamics}
\label{sec:frozen-dynamics}

We first isolate the explosion mechanism. Let \(Q=(Q_k)_{k\geq1}\) be the
tail of a \(\N\)-valued random variable \(Y\), and let \(Z\) be the
minimal increasing process which, from \(x<\infty\), updates at rate
\(x+1\) and moves to
\[
        \max\{x,U_x+Y'\},
\]
where \(U_x\) is uniform on \(\{0,\ldots,x\}\) and \(Y'\) is an independent
copy of \(Y\). Write \(\zeta:=\inf\{t:Z_t=\infty\}\).
Here, the minimal process means the usual pure-jump process
constructed up to its explosion time from the stated state-dependent
holding rates and independent graft marks, with \(\infty\) made absorbing.

\begin{proposition}[Frozen heavy-tail explosion]
\label{prop:frozen-heavy-tail}
Suppose that \(Q_k\geq ck^{-\alpha}\) for some \(c>0\) and
\(\alpha<1\). If \(Z_0\) has the law of \(Y\), then
\[
        \mathbb{P}(\zeta\leq s)>0,\qquad s>0.
\]
\end{proposition}

\begin{proof}
Fix \(x\geq1\). While \(x\leq Z_t<2x\), proposed updates whose graft
satisfies \(Y'\geq2x\) occur at rate at least \(xQ_{2x}\), and every such
update reaches level \(2x\). Uniformly in \(z\in[x,2x)\),
\begin{equation}
\label{eq:dyadic-crossing}
        \mathbb{E}_z[\tau_{2x}]
        \leq\frac1{xQ_{2x}}
        \leq Cx^{\alpha-1},
        \qquad
        \tau_{2x}:=\inf\{t:Z_t\geq2x\}.
\end{equation}

For \(K\geq1\) and \(z\geq K\), choose \(j\geq0\) such that
\(2^jK\leq z<2^{j+1}K\). Applying \eqref{eq:dyadic-crossing} first on
\([2^jK,2^{j+1}K)\) and then successively at the higher dyadic scales, the
strong Markov property gives
\[
        \mathbb{E}_z[\zeta]
        \leq
        C\sum_{i\geq j}(2^iK)^{\alpha-1}
        \leq C'K^{\alpha-1},
\]
uniformly in \(z\geq K\).
The series converges because \(\alpha<1\). 
Given \(s>0\), choose \(K\) so large that
\(C'K^{\alpha-1}\leq s/2\). Markov's inequality yields
\[
        \inf_{z\geq K}\mathbb{P}_z(\zeta\leq s)\geq\frac12,
\]
and therefore, since $Z_0$ is distributed as $Y$,
\[
        \mathbb{P}(\zeta\leq s)
        \geq
        \mathbb{P}(Y\geq K)\inf_{z\geq K}\mathbb{P}_z(\zeta\leq s)
        \geq\frac12Q_K>0.
\]
\end{proof}

Let \(Q_k:=q_k(t_*)\). By
Proposition~\ref{prop:critical-no-infinity}, \(Q_k\downarrow0\), so there is
a \(\N\)-valued random variable \(Y\) with tail
\(\mathbb{P}(Y\geq k)=Q_k\). Let \(Z\) be the corresponding frozen process,
started from \(Y\), and write \(r_k(s):=\mathbb{P}(Z_s\geq k)\).

\begin{proposition}[Postcritical domination]
\label{prop:postcritical-domination}
For every \(k\geq1\) and \(s\geq0\),
\[
        q_k(t_*+s)\geq r_k(s).
\]
Consequently,
\[
        q_\infty(t_*+s)\geq\mathbb{P}(\zeta\leq s).
\]
\end{proposition}

\begin{proof}
The frozen tails satisfy, for \(k\geq2\),
\begin{equation}
\label{eq:frozen-tail-system}
        r_k'
        =
        \sum_{u=1}^{k-1}Q_u(r_{k-u}-r_k)
        +Q_k(1-r_k),
        \qquad r_k(0)=Q_k.
\end{equation}
Indeed, this is the same crossing calculation as in
Proposition~\ref{prop:tail-system}, with the graft law frozen
at \(Q\).

Set \(\widetilde q_k(s):=q_k(t_*+s)\). Since
\(\widetilde q_u(s)\geq Q_u\),
\(\widetilde q_{k-u}-\widetilde q_k\geq0\), and
\(1-\widetilde q_k\geq0\), equation~\eqref{eq:tail-system} gives
\[
        \widetilde q_k'
        \geq
        \sum_{u=1}^{k-1}Q_u
        (\widetilde q_{k-u}-\widetilde q_k)
        +Q_k(1-\widetilde q_k).
\]
We prove \(\widetilde q_k\geq r_k\) by induction on \(k\). It is immediate
for \(k=1\). Assuming it holds below level \(k\), put
\(d_k:=\widetilde q_k-r_k\). Subtracting
\eqref{eq:frozen-tail-system},
\[
        d_k'
        +
        \left(\sum_{u=1}^{k-1}Q_u+Q_k\right)d_k
        \geq
        \sum_{u=1}^{k-1}Q_ud_{k-u}
        \geq0,
        \qquad d_k(0)=0.
\]
Integrating this inequality gives \(d_k\geq0\). Letting \(k\to\infty\) proves
the second assertion.
\end{proof}

\begin{proposition}[Equality of one-sided critical thresholds]
\label{prop:equality-critical-times}
One has \(t_{\mathrm{sus}}^{\mathrm{os}}=t_{\infty}^{\mathrm{os}}\), and
\[
        \mathbb{P}(X_{t_*+s}=\infty)>0,\qquad s>0.
\]
\end{proposition}


\begin{proof}
We already know that
\(t_{\mathrm{sus}}^{\mathrm{os}}\leq t_{\infty}^{\mathrm{os}}\).
By Proposition~\ref{prop:critical-no-infinity}, $X_{t_*}<\infty$ almost
surely, while Proposition~\ref{prop:critical-lower-tail} gives
\(Q_k\geq ck^{-2/3}\). Hence Propositions~\ref{prop:frozen-heavy-tail}
and~\ref{prop:postcritical-domination}, with \(\alpha=2/3\), imply that
for every \(s>0\),
\[
        q_\infty(t_*+s)\geq\mathbb{P}(\zeta\leq s)>0.
\]
Thus \(t_{\infty}^{\mathrm{os}}\leq t_*+s\) for every \(s>0\), hence
\(t_{\infty}^{\mathrm{os}}\leq t_*=t_{\mathrm{sus}}^{\mathrm{os}}\). The same estimate gives the
strictly postcritical explosion assertion.
\end{proof}

\begin{remark}
\label{rem:tail-explosion-criterion}
The frozen argument only requires a lower tail \(Q_k\geq ck^{-\alpha}\)
with \(\alpha<1\). The exponent \(2/3\) enters through the compensation
estimate of Proposition~\ref{prop:moment-comparison}, not through the
explosion mechanism.
\end{remark}

For the isotropic comparison we only need to restart the tail equation,
rather than the nonlinear process itself. The following elementary
observation records that the triangular system preserves the class of
tail sequences.

\begin{lemma}[Tail flow from general initial data]
\label{lem:tail-flow-general-data}
Let $\mu$ be a probability law on $\N$. There exists a unique family
$q^\mu=(q_k^\mu)_{k\geq1}$ such that $q_1^\mu\equiv1$,
$q_k^\mu\in C^1([0,\infty))$ for $k\geq2$,
$q_k^\mu(0)=\mu([k,\infty))$, and
\begin{equation}
\label{eq:general-tail-flow}
        (q_k^\mu)'
        =
        \sum_{u=1}^{k-1}q_u^\mu
        (q_{k-u}^\mu-q_k^\mu)
        +q_k^\mu(1-q_k^\mu),
        \qquad k\geq2.
\end{equation}
For every $t\geq0$,
\[
        1=q_1^\mu(t)\geq q_2^\mu(t)\geq\cdots\geq0,
\]
and each $q_k^\mu$ is non-decreasing in $t$. Consequently,
$q^\mu(t)$ is the tail of a unique probability law $\mu_t$ on
$\N\cup\{\infty\}$, with
\[
        \mu_t(\{\infty\})
        =
        q_\infty^\mu(t)
        :=
        \lim_{k\to\infty}q_k^\mu(t),
\]
and $(\mu_t)_{t\geq0}$ is non-decreasing in stochastic order.
\end{lemma}

\begin{proof}
Construct the coordinates recursively. Once
$q_1^\mu,\ldots,q_{k-1}^\mu$ are known, the equation for $q_k^\mu$ is the
scalar equation $x'=F_k(t,x)$, where
\[
        F_k(t,x)
        :=
        \sum_{u=1}^{k-1}q_u^\mu(t)
        (q_{k-u}^\mu(t)-x)+x(1-x).
\]
It is locally Lipschitz in $x$, while
$F_k(t,0)\geq0$ and $F_k(t,1)\leq0$. Hence $[0,1]$ is invariant, the
solution is global, and uniqueness follows recursively in $k$.

It remains to prove monotonicity in the level. We suppress the superscript
$\mu$ and set $D_k:=q_k-q_{k+1}$. Since $0\leq q_2\leq1$, one has
$D_1\geq0$. Suppose inductively that $D_j\geq0$ for $j<k$. Subtracting
the equations at levels $k$ and $k+1$ gives
\[
        D_k'
        =
        \sum_{u=1}^{k-1}q_uD_{k-u}
        -q_k
        -D_k\sum_{u=2}^{k}q_u
        +q_{k+1}^2.
\]
The induction hypothesis gives \(q_u\geq q_k\) for \(u<k\) and
\(\sum_{j=1}^{k-1}D_j=1-q_k\). Therefore
\[
\begin{aligned}
D_k'
&\geq
q_k(1-q_k)-q_k
-D_k\sum_{u=2}^{k}q_u+q_{k+1}^2 \\
&=
-D_k\left(
\sum_{u=2}^{k}q_u+q_k+q_{k+1}\right).
\end{aligned}
\]
Moreover,
\(D_k(0)=\mu(\{k\})\geq0\). An integrating factor therefore yields
\(D_k(t)\geq0\) for all \(t\), completing the induction.

Thus \(q_{k-u}\geq q_k\) for \(1\leq u<k\). Since also
\(0\leq q_u,q_k\leq1\), every term on the right-hand side of
\eqref{eq:general-tail-flow} is non-negative, and hence each
\(q_k^\mu\) is non-decreasing in \(t\).


Finally, setting
$\mu_t(\{k\})=q_k^\mu(t)-q_{k+1}^\mu(t)$ and
$\mu_t(\{\infty\})=q_\infty^\mu(t)$ defines a probability law by
telescoping, and its tail is $q^\mu(t)$. The last assertion follows from
the monotonicity of the tails in $t$.
\end{proof}

\begin{proposition}[Restarted one-sided sharpness]
\label{prop:restarted-one-sided}
Let $\mu$ be a probability law with finite support in $\N$, and let
$q^\mu$ be the tail flow of Lemma~\ref{lem:tail-flow-general-data}. Set
\[
        m_\mu(t):=\sum_{k\geq1}q_k^\mu(t),
        \qquad
        \tau_\mu:=\sup\{t\geq0:m_\mu(t)<\infty\}.
\]
Then
\begin{equation}
\label{eq:restarted-susceptibility-bound}
        \tau_\mu
        \leq
        2\log\left(1+\frac1{m_\mu(0)}\right),
\end{equation}
and
\begin{equation}
\label{eq:restarted-postcritical-explosion}
        q_\infty^\mu(\tau_\mu+s)>0,
        \qquad s>0.
\end{equation}
\end{proposition}

\begin{proof}
Since $\mu$ has finite support, all its moments are finite. On every
compact subinterval of $[0,\tau_\mu)$, the arguments of
Propositions~\ref{prop:first-moment} and~\ref{prop:polynomial-moments}
apply verbatim to \eqref{eq:general-tail-flow}. In particular, with
$A_\mu(t):=\sum_{k\geq1}kq_k^\mu(t)$,
\[
        m_\mu'
        =
        \frac12m_\mu^2+m_\mu
        -\frac12\sum_{k\geq1}(q_k^\mu)^2
        \geq\frac12m_\mu(m_\mu+1),
        \qquad
        A_\mu<\infty.
\]
Comparison with the solution of $y'=y(y+1)/2$ started from
$y(0)=m_\mu(0)$ proves
\eqref{eq:restarted-susceptibility-bound}. The truncated estimate
$m_{\mu,n}'\leq m_{\mu,n}^2$, with
$m_{\mu,n}:=\sum_{k=1}^nq_k^\mu$, gives as in
Proposition~\ref{prop:first-moment}
\[
        m_\mu(t)\longrightarrow\infty
        \qquad\text{as }t\uparrow\tau_\mu.
\]

Fix $t<\tau_\mu$ and normalize by
$p_k^\mu:=q_k^\mu/m_\mu$. Then $(p_k^\mu)_{k\geq1}$ is a non-increasing
probability mass function. Writing
\[
        L_\mu:=\sum_{1\leq u<k}kq_u^\mu q_k^\mu,
        \qquad
        D_\mu:=\sum_{k\geq1}k(q_k^\mu)^2,
\]
the same summation as in Proposition~\ref{prop:moment-comparison} gives
\[
        A_\mu'=2m_\mu A_\mu-L_\mu+A_\mu-D_\mu.
\]
Lemma~\ref{lem:decreasing-mass} yields
\[
        L_\mu
        \geq
        \frac23m_\mu A_\mu-\frac16m_\mu^2-\frac12D_\mu,
\]
and therefore
\[
        \frac{A_\mu'}{A_\mu}
        \leq\frac43(m_\mu+1)
        \leq\frac83\frac{m_\mu'}{m_\mu}.
\]
Hence
\[
        A_\mu(t)
        \leq
        A_\mu(0)
        \left(\frac{m_\mu(t)}{m_\mu(0)}\right)^{8/3}.
\]

For $t<\tau_\mu$, let $Y_t$ have law $\mu_t$. Then
$\mathbb E[Y_t]=m_\mu(t)$ and
$\mathbb E[Y_t^2]=2A_\mu(t)-m_\mu(t)$. Paley--Zygmund consequently gives
a constant $c_\mu>0$ such that
\[
        q_{\lceil m_\mu(t)/2\rceil}^\mu(t)
        \geq c_\mu m_\mu(t)^{-2/3}.
\]
For all sufficiently large integers $n$, choose
$t_n<\tau_\mu$ with $m_\mu(t_n)=2n$. Monotonicity in time gives
\[
        q_n^\mu(\tau_\mu)
        \geq q_n^\mu(t_n)
        \geq c_\mu' n^{-2/3}.
\]
The triangular system gives $q_k^\mu(t)>0$ for every $k$ and $t>0$;
after decreasing $c_\mu'$ over finitely many indices, the preceding
bound may therefore be taken for all $k\geq1$.

If $q_\infty^\mu(\tau_\mu)>0$, then
\eqref{eq:restarted-postcritical-explosion} follows immediately from
monotonicity in time. Otherwise
$Q_k:=q_k^\mu(\tau_\mu)\downarrow0$ is the tail of an almost surely finite
random variable and satisfies $Q_k\geq c_\mu'k^{-2/3}$.
Proposition~\ref{prop:frozen-heavy-tail} applies to the corresponding
frozen process. The comparison argument of
Proposition~\ref{prop:postcritical-domination}, with
$\widetilde q_k(s):=q_k^\mu(\tau_\mu+s)$, gives
\[
        q_\infty^\mu(\tau_\mu+s)
        \geq\mathbb P(\zeta\leq s)>0,
        \qquad s>0.
\]
\end{proof}

\section{Subcritical exponential moments}
\label{sec:exponential-moments}
Remember $(X_t)_{t\geq 0}$ stands for the shifted range process started from $1$.
\begin{proposition}
\label{prop:exponential-moments}
For every \(t_0<t_*\), there exists \(\lambda_{t_0}>0\) such that
\[
        \sup_{0\leq t\leq t_0}\mathbb{E}[e^{\lambda_{t_0}X_t}]<\infty.
\]
If \(t_0>0\), one may take
\begin{equation}
\label{eq:explicit-lambda}
        \lambda_{t_0}
        =
        \log\left(1+\frac{e^{-2t_0m(t_0)}}{2t_0}\right).
\end{equation}
\end{proposition}

\begin{proof}
Fix \(r>1\), and set
\[
        D_n(t):=\sum_{k=1}^n(r^k-1)q_k(t).
\]
By \eqref{eq:q-prime-gain-bound},
\[
\begin{aligned}
        D_n'(t)
        &\leq
        \sum_{\substack{u,v\geq1\\u+v\leq n}}
        (r^{u+v}-1)q_u(t)q_v(t)\\
        &\leq
        \sum_{u,v=1}^n(r^{u+v}-1)q_u(t)q_v(t).
\end{aligned}
\]
Since
\(r^{u+v}-1=(r^u-1)+(r^v-1)+(r^u-1)(r^v-1)\),
\begin{equation}
\label{eq:exponential-riccati}
        D_n'\leq2mD_n+D_n^2.
\end{equation}
Set
\[
        B(t):=\exp\left\{2\int_0^tm(s)\,ds\right\},
        \qquad
        H(t):=\int_0^tB(s)\,ds.
\]
Writing \(r=1+\delta\), one has \(D_n(0)=\delta\). For
\(E_n:=D_n/B\), inequality~\eqref{eq:exponential-riccati} gives
\(E_n'\leq BE_n^2\), hence
\[
        \frac1{E_n(t)}\geq\frac1\delta-H(t).
\]
Whenever \(\delta H(t_0)<1\),
\[
        D_n(t)
        \leq
        \frac{\delta B(t)}{1-\delta H(t)},
        \qquad 0\leq t\leq t_0.
\]
Since \(m\) is non-decreasing,
\(B(t)\leq e^{2t_0m(t_0)}\) and
\(H(t_0)\leq t_0e^{2t_0m(t_0)}\). For \(t_0>0\), choose
\(\delta=e^{-2t_0m(t_0)}/(2t_0)\). Then
\(\delta H(t_0)\leq1/2\), and the preceding estimate is uniform in \(n\)
and \(0\leq t\leq t_0\). Letting \(n\to\infty\),
\[
        \sup_{0\leq t\leq t_0}
        \sum_{k\geq1}(r^k-1)q_k(t)<\infty.
\]
For every \(\N\)-valued random variable \(W\), and $r\geq 1$,
\begin{equation}
\label{eq:generic-W}
        \mathbb{E}[r^W]-1
        =
        (r-1)\sum_{k\geq1}r^{k-1}\mathbb{P}(W\geq k)  \leq
        \sum_{k\geq1}(r^k-1) \mathbb{P}(W\geq k).
\end{equation}
Applying this identity to \(X_t\) proves the assertion with
\(\lambda_{t_0}=\log r\). The case \(t_0=0\) is immediate.
\end{proof}

\begin{corollary}
\label{cor:subcritical-tails}
For every \(t<t_*\), there exist \(c_t,C_t>0\) such that
\[
        \mathbb{P}(X_t\geq k)\leq C_te^{-c_tk},
        \qquad k\geq1.
\]
In particular, \(\mathbb{E}[X_t^p]<\infty\) for every \(p>0\).
\end{corollary}

\begin{proof}
Choose \(t_0\in(t,t_*)\) and apply Markov's inequality with the exponential
moment furnished by Proposition~\ref{prop:exponential-moments}.
\end{proof}

\begin{proof}[Proof of Theorem~\ref{thm:range-criticality}]
The bounds on $t_{\mathrm{sus}}^{\mathrm{os}}$ and the divergence of
$m$ follow from Proposition~\ref{prop:first-moment}. The critical lower tail
is Proposition~\ref{prop:critical-lower-tail}, while
Proposition~\ref{prop:critical-no-infinity} shows that
$X_{t_*}<\infty$ almost surely. Proposition~\ref{prop:equality-critical-times}
identifies $t_{\mathrm{sus}}^{\mathrm{os}}$ with $t_{\infty}^{\mathrm{os}}$
and gives strictly postcritical explosion. The subcritical exponential
moments follow from Proposition~\ref{prop:exponential-moments}.
\end{proof}

\section{Endpoint domination and isotropic sharpness}
\label{sec:isotropic-range}

We now consider the one-dimensional isotropic Bak--Sneppen model, in which
the minimum site and both of its neighbours are refreshed. Starting from a
configuration above level $t$, force one update of $\{-1,0,1\}$ and continue
until all fitnesses again exceed $t$. Its infinite-volume range is an interval
\[
        \mathcal R_t=[-L_t,R_t],
\]
where $L_t,R_t\in\N\cup\{\infty\}$ and $L_0=R_0=1$. Put
\[
        S_t:=L_t+R_t+1,
\]
with the convention that $S_t=\infty$ as soon as one endpoint is infinite,
and define
\[
        t_{\mathrm{sus}}^{\mathrm{iso}}
        :=\sup\{t\geq0:\mathbb{E}[S_t]<\infty\},
        \qquad
        t_{\infty}^{\mathrm{iso}}
        :=\inf\{t\geq0:\mathbb{P}(S_t=\infty)>0\}.
\]

Positive mass at infinity forces an infinite mean, so
\begin{equation}
\label{eq:isotropic-easy-threshold-order}
        t_{\mathrm{sus}}^{\mathrm{iso}}
        \leq t_{\infty}^{\mathrm{iso}}.
\end{equation}
The level-regeneration and graphical arguments of
Section~\ref{sec:avalanche-range} apply without change, except that the
position of the triggering site inside the current interval must be retained, and the random directed graph is defined on $\Z \times [0,\infty)$.

\begin{proposition}[Isotropic range evolution]
\label{prop:isotropic-range-evolution}
At level $t$, conditionally on
$(L_t,R_t)=(\ell,r)\in\N^2$, the proposed graft rate is
$\ell+r+1$. If $U$ is uniform on $\{-\ell,\ldots,r\}$ and
$(L_t',R_t')$ is an independent copy of $(L_t,R_t)$, then the proposed
update is
\begin{equation}
\label{eq:isotropic-range-update}
        (\ell,r)
        \longmapsto
        \bigl(\max\{\ell,L_t'-U\},\max\{r,R_t'+U\}\bigr).
\end{equation}
Null updates are allowed.
\end{proposition}

\begin{proof}
The current range contains $\ell+r+1$ sites, each carrying a rate-one
Poisson process on the level axis. Their superposition has rate
$\ell+r+1$, and the location $U$ of the next mark is uniform on the current
interval. The mark carries an independent isotropic avalanche at the current
level. After translation by $U$, its range is
$[U-L_t',U+R_t']$; taking the union with $[-\ell,r]$ gives
\eqref{eq:isotropic-range-update}.
\end{proof}


For $k\geq1$, set
\[
h_k(t):=\mathbb{P}(L_t\geq k),
\qquad
\ell(t):=\mathbb{E}[L_t]=\sum_{k\geq1}h_k(t),
\]
where the equality is understood in $[0,\infty]$. Reflection symmetry gives
$L_t\stackrel{\mathrm d}=R_t$ and
$\mathbb{E}[S_t]=2\ell(t)+1$. For a tail sequence $a=(a_k)_{k\geq1}$ with
$a_1=1$, define
\[
\Phi_k(a)
:=
\sum_{u=1}^{k-1}a_u(a_{k-u}-a_k)+a_k(1-a_k),
\qquad k\geq2.
\]
\begin{proposition}[Endpoint tail inequalities]
\label{prop:isotropic-endpoint-tail}
For every $0\leq s\leq t$ and $k\geq2$,
\begin{equation}
\label{eq:isotropic-endpoint-supersolution}
h_k(t)-h_k(s)
\geq
\int_s^t\Phi_k(h(u))\,du.
\end{equation}
Moreover, $h_k$ is continuously differentiable on
$[0,t_{\mathrm{sus}}^{\mathrm{iso}})$ and, for every
$t<t_{\mathrm{sus}}^{\mathrm{iso}}$,
\begin{equation}
\label{eq:isotropic-endpoint-upper}
h_k'(t)
\leq
\sum_{u=1}^{k-1}h_u(t)h_{k-u}(t)
+\ell(t)h_k(t).
\end{equation}
\end{proposition}
\begin{proof}
Fix $k\geq2$. For $0\leq s\leq t$, let $N_k(s,t]$ be the number of
marks $(-v,u)$ with $s<u\leq t$ such that $L_{u-}<k$,
$0\leq v\leq L_{u-}$, and the attached avalanche has left endpoint at
least $k-v$. Any such mark makes the left endpoint reach $k$, and no
further such mark is counted thereafter. Hence
\[
N_k(s,t]
\leq
\mathbf{1}_{\{L_t\geq k\}}
-\mathbf{1}_{\{L_s\geq k\}}.
\]
Only the $k$ fixed vertical lines $\{-(k-1),\ldots,0\}$ are involved, so
compensation gives
\[
\begin{aligned}
h_k(t)-h_k(s)
&\geq
\int_s^t
\mathbb{E}\left[
\mathbf{1}_{\{L_u<k\}}
\sum_{v=0}^{L_u}h_{k-v}(u)
\right]du  \\
&=
\int_s^t\sum_{\ell=1}^{k-1}
(h_\ell(u)-h_{\ell+1}(u))
\sum_{v=0}^{\ell}h_{k-v}(u)\,du
=
\int_s^t\Phi_k(h(u))\,du,
\end{aligned}
\]
where the last equality follows by the same summation by parts as in
Proposition~\ref{prop:tail-system}. This proves
\eqref{eq:isotropic-endpoint-supersolution}.
Fix now $t_0<t_{\mathrm{sus}}^{\mathrm{iso}}$. Since the range is
increasing,
\[
\int_0^{t_0}\mathbb{E}[S_u]\,du
\leq t_0\mathbb{E}[S_{t_0}]<\infty,
\]
so the full crossing process has an integrable compensator on
$[0,t_0]$. Conditionally on $(L_u,R_u)=(\ell,r)$ with $\ell<k$, its
rate of reaching level $k$ is
$\sum_{v=0}^{\ell}h_{k-v}(u)+\sum_{j=1}^{r}h_{k+j}(u)$. Therefore, for
$0\leq s\leq t\leq t_0$,
\[
h_k(t)-h_k(s)
=
\int_s^t\left\{
\Phi_k(h(u))
+
\mathbb{E}\left[
\mathbf{1}_{\{L_u<k\}}
\sum_{j=1}^{R_u}h_{k+j}(u)
\right]\right\}du.
\]
This shows that $h_k$ is absolutely
continuous on $[0,t_0]$,  
and the same is true of $h_1$ since $h_1\equiv1$.
We claim the integrand in the preceding identity is continuous. Indeed, the
number of reachable marks on $[0,t_0]$ has finite expectation and is
therefore almost surely finite, while a fixed deterministic level is
almost surely not a mark. Thus $(L_v,R_v)\to(L_u,R_u)$ almost surely as
$v\to u$. Since every $h_j$ is continuous and
\[
\mathbf{1}_{\{L_v<k\}}
\sum_{j=1}^{R_v}h_{k+j}(v)
\leq R_{t_0},
\qquad \mathbb{E}[R_{t_0}]<\infty,
\]
dominated convergence proves the claim. Hence $h_k$ is continuously
differentiable and its derivative equals the displayed integrand.
Finally, using $h_1=1$, monotonicity of the tails, and reflection
symmetry,
\[
\Phi_k(h(t))
\leq
\sum_{u=1}^{k-1}h_u(t)h_{k-u}(t),
\]
while
\[
\mathbb{E}\left[
\mathbf{1}_{\{L_t<k\}}
\sum_{j=1}^{R_t}h_{k+j}(t)
\right]
\leq
\mathbb{E}[R_t]h_k(t)
=
\ell(t)h_k(t).
\]
This proves \eqref{eq:isotropic-endpoint-upper}. Since $t_0$ was
arbitrary, the result holds throughout
$[0,t_{\mathrm{sus}}^{\mathrm{iso}})$.
\end{proof}
\begin{remark}
A direct repetition of the Paley--Zygmund argument for the isotropic
endpoint loses too much. Indeed, set
$A_L(t):=\sum_{k\geq1}k h_k(t)$. In the exact decomposition of the proof
above, the $\Phi_k$ term is controlled by the same compensation argument
as Proposition~\ref{prop:moment-comparison}, contributing at most
$\frac43(\ell+1)A_L$, whereas the additional contribution from grafts
rooted to the right is at most $\ell A_L$. Hence
\[
A_L'\leq\left(\frac73\ell+\frac43\right)A_L,
\qquad
\frac{\ell'}{\ell}\geq\frac{\ell+1}{2},
\]
and therefore $A_L\lesssim\ell^{14/3}$. Paley--Zygmund then yields only
$\mathbb P(L\geq k)\gtrsim k^{-8/3}$ at criticality, too weak for
Proposition~\ref{prop:frozen-heavy-tail}. This motivates the endpoint
comparison below.
\end{remark}
We write $\preceq$ for the usual stochastic order.
\begin{proposition}[One-sided endpoint domination]
\label{prop:isotropic-endpoint-domination}
Fix $s\geq0$ with $\mathbb E[L_s]<\infty$ and $K\geq1$. Let
$q^{s,K}$ be the tail flow of Lemma~\ref{lem:tail-flow-general-data}
with initial data
\[
q_k^{s,K}(0)
:=
\mathbb P(L_s\wedge K\geq k).
\]
Then, for every $k\geq1$ and $t\geq0$,
\begin{equation}
\label{eq:isotropic-endpoint-domination}
q_k^{s,K}(t)
\leq
\mathbb P(L_{s+t}\geq k).
\end{equation}
Consequently, with $m_{s,K}:=\mathbb E[L_s\wedge K]$,
\begin{equation}
\label{eq:isotropic-explosion-upper}
t_{\infty}^{\mathrm{iso}}
\leq
s+2\log\left(1+\frac1{m_{s,K}}\right).
\end{equation}
In particular,
\begin{equation}
\label{eq:isotropic-one-sided-threshold-comparison}
X_t\preceq L_t,
\qquad
t_{\infty}^{\mathrm{iso}}
\leq t_{\infty}^{\mathrm{os}}.
\end{equation}
\end{proposition}
\begin{proof}
Set $\widetilde h_k(t):=h_k(s+t)$. Then
$(q_k^{s,K})'=\Phi_k(q^{s,K})$,
$q_k^{s,K}(0)\leq\widetilde h_k(0)$, and
\eqref{eq:isotropic-endpoint-supersolution} shows that
$\widetilde h_k$ is an integral supersolution of the same equation.
We prove $q_k^{s,K}\leq\widetilde h_k$ by induction on $k$. The assertion
is immediate for $k=1$.
Suppose it holds below level $k$ and set
$d_k:=q_k^{s,K}-\widetilde h_k$. For $0\leq r\leq t$,
\[
d_k(t)-d_k(r)
\leq
\int_r^t
\bigl(\Phi_k(q^{s,K}(u))-\Phi_k(\widetilde h(u))\bigr)\,du.
\]

For fixed \(x\in[0,1]\), consider
\[
        \Psi_x(a_1,\ldots,a_{k-1})
        :=
        \sum_{u=1}^{k-1}a_u(a_{k-u}-x)+x(1-x).
\]
For \(j<k\),
\[
        \partial_{a_j}\Psi_x
        =2a_{k-j}-x.
\]
Hence \(\Psi_x\) is non-decreasing in each coordinate on the region
where \(a_i\geq x\) for every \(i<k\).

By the induction hypothesis,
\(q^{s,K}_j(u)\leq\widetilde h_j(u)\) for \(j<k\), while
Lemma~\ref{lem:tail-flow-general-data} gives
\(q^{s,K}_j(u)\geq q^{s,K}_k(u)\). Therefore, with the \(k\)-th
coordinate fixed at \(q^{s,K}_k(u)\), increasing the lower coordinates
from \(q^{s,K}(u)\) to \(\widetilde h(u)\) can only increase
\(\Phi_k\). Consequently,
\[
\begin{aligned}
        \Phi_k(q^{s,K}(u))-\Phi_k(\widetilde h(u))
        &\leq
        \Phi_k\bigl(
        \widetilde h_1(u),\ldots,\widetilde h_{k-1}(u),
        q^{s,K}_k(u)
        \bigr)
        -\Phi_k(\widetilde h(u))  \\
        &=
        a_k(u)d_k(u),
\end{aligned}
\]
where
\[
        a_k(u)
        :=
        1-\sum_{j=1}^{k-1}\widetilde h_j(u)
        -q^{s,K}_k(u)-\widetilde h_k(u)
        \leq0,
\]
using $\tilde h_1 \equiv 1$. 
Thus
\begin{equation}
\label{eq:isotropic-endpoint-comparison-difference}
d_k(t)-d_k(r)
\leq
\int_r^t a_k(u)d_k(u)\,du.
\end{equation}
No continuity of $\widetilde h_k$ is needed here. Since
$\widetilde h_k$ is non-decreasing, every discontinuity of
$d_k=q_k^{s,K}-\widetilde h_k$ is downward; the integral inequality
\eqref{eq:isotropic-endpoint-comparison-difference} controls any possible
increase.
We claim that $d_k\leq0$. Otherwise, fix $t>0$ such that $d_k(t)>0$ and
let
\[
\tau:=\sup\{r\in[0,t]:d_k(r)\leq0\}.
\]
The set is nonempty since $d_k(0)\leq0$. On $(\tau,t]$ one has
$d_k>0$, and hence $a_kd_k\leq0$. Choose $r_n\uparrow\tau$ with
$d_k(r_n)\leq0$; if $\tau=0$, take $r_n=0$. Since $a_k$ and $d_k$ are
bounded on compact intervals, \eqref{eq:isotropic-endpoint-comparison-difference}
gives
\[
\begin{aligned}
d_k(t)
&\leq
d_k(r_n)
+\int_{r_n}^{\tau}|a_k(u)d_k(u)|\,du
+\int_{\tau}^{t}a_k(u)d_k(u)\,du \\
&\leq
\int_{r_n}^{\tau}|a_k(u)d_k(u)|\,du
\longrightarrow0.
\end{aligned}
\]
If $\tau=t$, the same estimate applies with the last integral absent.
In either case this contradicts $d_k(t)>0$. Thus $d_k\leq0$, completing
the induction and proving \eqref{eq:isotropic-endpoint-domination}.
Set
\[
\tau_{s,K}
:=
\sup\left\{t\geq0:
\sum_{k\geq1}q_k^{s,K}(t)<\infty\right\}.
\]
Proposition~\ref{prop:restarted-one-sided} gives
\[
\tau_{s,K}
\leq
2\log\left(1+\frac1{m_{s,K}}\right),
\qquad
q_\infty^{s,K}(\tau_{s,K}+\varepsilon)>0
\]
for every $\varepsilon>0$. Letting $k\to\infty$ in
\eqref{eq:isotropic-endpoint-domination} at time
$\tau_{s,K}+\varepsilon$ yields
\[
\mathbb P
(L_{s+\tau_{s,K}+\varepsilon}=\infty)
\geq
q_\infty^{s,K}(\tau_{s,K}+\varepsilon)>0.
\]
Hence
\[
t_{\infty}^{\mathrm{iso}}
\leq
s+2\log\left(1+\frac1{m_{s,K}}\right)+\varepsilon,
\]
and \eqref{eq:isotropic-explosion-upper} follows by letting
$\varepsilon\downarrow0$.
For $s=0$ and $K=1$, the initial data are
$q_1^{0,1}(0)=1$ and $q_k^{0,1}(0)=0$ for $k\geq2$. Uniqueness of the solution to the
triangular system therefore gives
$q_k^{0,1}(t)=\mathbb P(X_t\geq k)$. Thus
\eqref{eq:isotropic-endpoint-domination} gives $X_t\preceq L_t$, and
letting $k\to\infty$ gives
$t_{\infty}^{\mathrm{iso}}\leq t_{\infty}^{\mathrm{os}}$.
\end{proof}



\begin{lemma}[Finite-mean continuation]
\label{lem:isotropic-finite-mean-continuation}
Suppose that $M(s):=\mathbb{E}[S_s]<\infty$. Then
\begin{equation}
\label{eq:isotropic-mean-continuation}
        \mathbb{E}[S_{s+t}]
        \leq
        \frac{M(s)}{1-tM(s)},
        \qquad 0\leq t<M(s)^{-1}.
\end{equation}
\end{lemma}

\begin{proof}
For $K\geq1$, set $M_K(t):=\mathbb{E}[S_t\wedge K]$. The
truncated observable can change only while $S_t<K$; there its jump rate
and jump size are bounded by $K$. Thus $M_K$ is locally absolutely
continuous. Conditionally on $S_t=x<K$, marks occur at rate $x$, and
adjoining an independent graft of size $S_t'$ increases the truncated
size by at most $S_t'\wedge K$. The compensator identity therefore
gives, for almost every $t$,
\[
        M_K'(t)
        \leq
        \mathbb{E}[S_t\mathbf 1_{\{S_t<K\}}]M_K(t)
        \leq M_K(t)^2.
\]
Integrating from $s$ to $s+t$ and using then
$M_K(s)\leq M(s)$ gives
\[
        M_K(s+t)
        \leq
        \frac{M(s)}{1-tM(s)},
        \qquad 0\leq t<M(s)^{-1}.
\]
Letting $K\to\infty$ and using monotone convergence proves
\eqref{eq:isotropic-mean-continuation}.
\end{proof}

\begin{theorem}[Isotropic sharpness]
\label{thm:isotropic-criticality}
In the isotropic model,
\begin{equation}
\label{eq:isotropic-criticality}
        t_{\mathrm{sus}}^{\mathrm{iso}}
        =t_{\infty}^{\mathrm{iso}}
        =:t_c^{\mathrm{iso}}.
\end{equation}
Moreover,
\begin{equation}
\label{eq:isotropic-critical-bounds}
        \frac13
        \leq t_c^{\mathrm{iso}}
        \leq t_c^{\mathrm{os}}
        \leq2\log2,
\end{equation}
and
\[
        \mathbb{P}(S_{t_c^{\mathrm{iso}}+s}=\infty)>0,
        \qquad s>0.
\]
\end{theorem}

\begin{proof}
Lemma~\ref{lem:isotropic-finite-mean-continuation} at $s=0$, where
$S_0=3$, gives $t_{\mathrm{sus}}^{\mathrm{iso}}\geq1/3$.
Proposition~\ref{prop:isotropic-endpoint-domination} and
Theorem~\ref{thm:range-criticality} give
$t_{\infty}^{\mathrm{iso}}\leq t_c^{\mathrm{os}}\leq2\log2$; hence
$t_{\mathrm{sus}}^{\mathrm{iso}}<\infty$.

We first claim that $\mathbb{E}[S_t]\to\infty$ as
$t\uparrow t_{\mathrm{sus}}^{\mathrm{iso}}$. Otherwise
$\mathbb{E}[S_t]\leq C$ for all $t<t_{\mathrm{sus}}^{\mathrm{iso}}$ and some
finite $C$. Choose $s<t_{\mathrm{sus}}^{\mathrm{iso}}$ so close to
$t_{\mathrm{sus}}^{\mathrm{iso}}$ that
$t_{\mathrm{sus}}^{\mathrm{iso}}-s<(2C)^{-1}$. Since
$\mathbb{E}[S_s]\leq C$, Lemma~\ref{lem:isotropic-finite-mean-continuation}
gives $\mathbb{E}[S_{s+(2C)^{-1}}]<\infty$, although
$s+(2C)^{-1}>t_{\mathrm{sus}}^{\mathrm{iso}}$, contradicting the definition
of the susceptibility threshold. By reflection symmetry,
$\mathbb{E}[L_t]=(\mathbb{E}[S_t]-1)/2\to\infty$ as
$t\uparrow t_{\mathrm{sus}}^{\mathrm{iso}}$.

Fix $t>t_{\mathrm{sus}}^{\mathrm{iso}}$. Choose
$s<t_{\mathrm{sus}}^{\mathrm{iso}}$ so that
$2\log(1+1/\mathbb{E}[L_s])<t-s$. Since
$\mathbb{E}[L_s\wedge K]\uparrow\mathbb{E}[L_s]$, the same strict inequality
holds for all sufficiently large $K$ with $\mathbb{E}[L_s\wedge K]$ in
place of $\mathbb{E}[L_s]$. Proposition~\ref{prop:isotropic-endpoint-domination}
and Proposition~\ref{prop:restarted-one-sided} then give
$\mathbb{P}(L_t=\infty)>0$, hence $\mathbb{P}(S_t=\infty)>0$. Therefore
$t_{\infty}^{\mathrm{iso}}\leq t_{\mathrm{sus}}^{\mathrm{iso}}$, which
together with \eqref{eq:isotropic-easy-threshold-order} proves
\eqref{eq:isotropic-criticality}. The bounds
\eqref{eq:isotropic-critical-bounds} and the strictly postcritical assertion
have already been established.
\end{proof}



The next two results parallel their one-sided counterparts,
Propositions~\ref{prop:critical-no-infinity}
and~\ref{prop:exponential-moments}. Their proofs adapt the same
arguments using the isotropic comparison and the endpoint estimates of
Proposition~\ref{prop:isotropic-endpoint-tail}.


\begin{proposition}[No infinite avalanche at criticality]
\label{prop:isotropic-critical-finiteness}
At the isotropic critical level,
\[
        \mathbb{P}(S_{t_c^{\mathrm{iso}}}=\infty)=0.
\]
\end{proposition}

\begin{proof}
Taking $s=t$ in \eqref{eq:isotropic-explosion-upper} and then letting
$K\to\infty$ gives, for $t<t_c^{\mathrm{iso}}$,
\[
\begin{aligned}
        t_c^{\mathrm{iso}}-t
        &\leq2\log\left(1+\frac1{\mathbb{E}[L_t]}\right),\\
        \mathbb{E}[S_t]
        =2\mathbb{E}[L_t]+1
        &\leq\frac{2}{e^{(t_c^{\mathrm{iso}}-t)/2}-1}+1
        \leq\frac{4}{t_c^{\mathrm{iso}}-t}+1.
\end{aligned}
\]
For $n\geq2$, set $\rho_n(t):=\mathbb{P}(S_t\geq n)$. Fix $C>0$ and, for
all sufficiently large $n$, let $t_n:=t_c^{\mathrm{iso}}-C/n$. Markov's
inequality gives
\[
        \rho_n(t_n)
        \leq\frac{\mathbb{E}[S_{t_n}]}n
        \leq\frac4C+\frac1n.
\]

Set
$b_n:=(n-1)^{-1}\sum_{j=1}^{n-1}\rho_j(t_c^{\mathrm{iso}})$. Before $S_t$
reaches $n$, proposed grafts occur at total rate $S_t$. Counting every
proposed graft as a possible crossing of level $n$ therefore bounds the
crossing intensity by
\[
        \mathbb E\left[S_u\mathbf 1_{\{S_u<n\}}\right]
        =\sum_{j=1}^{n-1}\rho_j(u)-(n-1)\rho_n(u).
\]
Consequently, the compensator identity and monotonicity in the level give,
for $t_n\leq s\leq r<t_c^{\mathrm{iso}}$,
\[
\begin{aligned}
        \rho_n(r)-\rho_n(s)
        &\leq
        \int_s^r\left(
        \sum_{j=1}^{n-1}\rho_j(u)-(n-1)\rho_n(u)\right)du \\
        &\leq
        (n-1)\int_s^r\bigl(b_n-\rho_n(u)\bigr)\,du.
\end{aligned}
\]
An integrating factor, followed by $r\uparrow t_c^{\mathrm{iso}}$, yields
\[
        \rho_n(t_c^{\mathrm{iso}})
        \leq
        e^{-(n-1)C/n}\rho_n(t_n)
        +b_n\bigl(1-e^{-(n-1)C/n}\bigr).
\]
Here $\rho_n(r)\uparrow\rho_n(t_c^{\mathrm{iso}})$ as
$r\uparrow t_c^{\mathrm{iso}}$: in the graphical construction, before the
range reaches $n$ at most $n-1$ vertical lines are active, and almost surely
none has a mark at the deterministic level $t_c^{\mathrm{iso}}$.
Moreover, both $\rho_n(t_c^{\mathrm{iso}})$ and $b_n$ converge to
$\mathbb{P}(S_{t_c^{\mathrm{iso}}}=\infty)$, the latter by Ces\`aro
convergence. Rearranging the preceding inequality and using
$\rho_n(t_n)\leq\rho_n(t_c^{\mathrm{iso}})$ therefore gives
\[
        \rho_n(t_n)\longrightarrow
        \mathbb{P}(S_{t_c^{\mathrm{iso}}}=\infty).
\]
The Markov bound above now gives
$\mathbb{P}(S_{t_c^{\mathrm{iso}}}=\infty)\leq4/C$. Since $C>0$ is
arbitrary, the probability is zero.
\end{proof}

\begin{proposition}[Subcritical isotropic exponential moments]
\label{prop:isotropic-exponential-moments}
For every $t_0<t_c^{\mathrm{iso}}$, there exists $\lambda_{t_0}>0$ such that
\[
        \sup_{0\leq t\leq t_0}
        \mathbb{E}[e^{\lambda_{t_0}S_t}]<\infty.
\]
\end{proposition}

\begin{proof}
Fix $r>1$ and, for $n\geq1$, set
\[
        D_n(t):=\sum_{k=1}^n(r^k-1)h_k(t).
\]
By Proposition~\ref{prop:isotropic-endpoint-tail}, for every
$t\leq t_0$,
\begin{align*}
        D_n'(t)
        &\leq
        \sum_{\substack{u,v\geq1\\u+v\leq n}}
        (r^{u+v}-1)h_u(t)h_v(t)+\ell(t)D_n(t)\\
        &\leq 3\ell(t)D_n(t)+D_n(t)^2.
\end{align*}
Indeed, the convolution term is bounded by
$2\ell(t)D_n(t)+D_n(t)^2$, exactly as in
\eqref{eq:exponential-riccati}. Set
\[
        B(t):=\exp\left\{3\int_0^t\ell(u)\,du\right\},
        \qquad
        H(t):=\int_0^tB(u)\,du.
\]
Writing $r=1+\delta$, one has $D_n(0)=\delta$. For
$E_n:=D_n/B$, the preceding inequality gives $E_n'\leq BE_n^2$ and hence
\[
        D_n(t)
        \leq
        \frac{\delta B(t)}{1-\delta H(t)},
        \qquad 0\leq t\leq t_0,
\]
whenever $\delta H(t_0)<1$. Since $t_0<t_c^{\mathrm{iso}}$,
$\ell(t_0)<\infty$ and therefore $H(t_0)<\infty$; choose
$\delta>0$ accordingly. Letting $n\to\infty$ and using  \eqref{eq:generic-W} with $L_t$ 
% \[
%         (r-1)\sum_{k\geq1}r^{k-1}h_k(t)
%         \leq
%         \sum_{k\geq1}(r^k-1)h_k(t)
% \]
% together with the tail-sum identity for $\mathbb{E}[r^{L_t}]$ 
gives
\[
        \sup_{0\leq t\leq t_0}\mathbb{E}[r^{L_t}]<\infty.
\]
With $\lambda_{t_0}:=\frac12\log r$, reflection symmetry and
Cauchy--Schwarz give
\[
        \mathbb{E}[e^{\lambda_{t_0}S_t}]
        =e^{\lambda_{t_0}}
        \mathbb{E}\left[r^{(L_t+R_t)/2}\right]
        \leq
        e^{\lambda_{t_0}}
        \bigl(\mathbb{E}[r^{L_t}]\mathbb{E}[r^{R_t}]\bigr)^{1/2},
\]
uniformly for $0\leq t\leq t_0$.
\end{proof}

\section{Stationary product limits}
\label{sec:stationary-limits}

Let $\Gamma:=\{\mathrm{os},\mathrm{iso}\}$, and write
\[
        \mathcal N_{\mathrm{os}}:=\{0,1\},
        \qquad
        \mathcal N_{\mathrm{iso}}:=\{-1,0,1\},
        \qquad
        d_\gamma:=|\mathcal N_\gamma|.
\]
For $N\geq2$, let $\Lambda_N:=\Z/N\Z$ and consider the Bak--Sneppen
chain in which the sites in $i+\mathcal N_\gamma$ are refreshed when $i$
is the minimum site, with repetitions modulo $N$ ignored. Let
$r_N^\gamma(t)$ and $\eta_N^\gamma(t)$ denote the range and duration of a
finite-volume $t$-avalanche, and let $r_\infty^\gamma(t)$ and
$\eta_\infty^\gamma(t)$ be their infinite-volume counterparts. Thus
\[
        r_\infty^{\mathrm{os}}(t)=X_t+1,
        \qquad
        r_\infty^{\mathrm{iso}}(t)=S_t.
\]
Set
\[
\begin{aligned}
        R_N^\gamma(t)&:=\mathbb{E}[r_N^\gamma(t)],
        &D_N^\gamma(t)&:=\mathbb{E}[\eta_N^\gamma(t)],
        &P_N^\gamma(t)&:=\mathbb{P}(r_N^\gamma(t)=N),\\
        R_\infty^\gamma(t)&:=\mathbb{E}[r_\infty^\gamma(t)],
        &D_\infty^\gamma(t)&:=\mathbb{E}[\eta_\infty^\gamma(t)],
        &P_\infty^\gamma(t)&:=\mathbb{P}(r_\infty^\gamma(t)=\infty),
\end{aligned}
\]
and recall
\[
\begin{aligned}
        t_{\mathrm d}^\gamma
        &:=\sup\{t\geq0:D_\infty^\gamma(t)<\infty\},\\
        t_{\mathrm{sus}}^\gamma
        &:=\sup\{t\geq0:R_\infty^\gamma(t)<\infty\},\\
        t_\infty^\gamma
        &:=\inf\{t\geq0:P_\infty^\gamma(t)>0\}.
\end{aligned}
\]
Every site in the avalanche range is refreshed during one of its updates,
so $r_\infty^\gamma(t)\leq d_\gamma\eta_\infty^\gamma(t)$. Hence
\[
        t_{\mathrm d}^\gamma
        \leq t_{\mathrm{sus}}^\gamma
        \leq t_\infty^\gamma.
\]
\begin{theorem}[Stationary product limit]
\label{thm:stationary-product-limits}
For each $\gamma\in\Gamma$,
\[
        t_{\mathrm d}^\gamma
        =t_{\mathrm{sus}}^\gamma
        =t_\infty^\gamma
        =t_c^\gamma,
\]
where $t_c^{\mathrm{os}}$ and $t_c^{\mathrm{iso}}$ are the critical values
of Theorems~\ref{thm:range-criticality} and
\ref{thm:isotropic-criticality}, respectively. If
\[
        F_N^\gamma
        =\bigl(F_{N,x}^\gamma\bigr)_{x\in\Lambda_N}
\]
has law $\pi_N^\gamma$, then, for every fixed collection of distinct sites
$x_1,\ldots,x_\ell\in\mathbb Z$,
\[
        \bigl(
        F_{N,x_1}^\gamma,\ldots,F_{N,x_\ell}^\gamma
        \bigr)
        \xrightarrow[N\to\infty]{\mathrm d}
        \nu_{t_c^\gamma}^{\otimes\ell}.
\]
Equivalently, in the original uniform parametrisation, the limiting
fitnesses are independent and uniform on $[q_c^\gamma,1]$, where $q_c^\gamma:=1-e^{-t_c^\gamma}$.
\end{theorem}

\begin{proof}
Theorem~\ref{thm:range-criticality} gives
\[
        t_{\mathrm{sus}}^{\mathrm{os}}
        =t_\infty^{\mathrm{os}}
        =t_c^{\mathrm{os}},
        \qquad
        1\leq t_c^{\mathrm{os}}\leq2\log2,
\]
while Theorem~\ref{thm:isotropic-criticality} gives
\[
        t_{\mathrm{sus}}^{\mathrm{iso}}
        =t_\infty^{\mathrm{iso}}
        =t_c^{\mathrm{iso}},
        \qquad
        \frac13\leq t_c^{\mathrm{iso}}
        \leq t_c^{\mathrm{os}}.
\]
Proposition~\ref{prop:app-range-duration} yields
$t_{\mathrm d}^\gamma=t_{\mathrm{sus}}^\gamma$, and
Proposition~\ref{prop:app-stationary-limit} gives the convergence of the
stationary finite-dimensional marginals. Finally, if
$F=t_c^\gamma+E$ with $E\sim\operatorname{Exp}(1)$, then $e^{-E}$ is
uniform on $[0,1]$, so $1-e^{-F}$ is uniform on
$[1-e^{-t_c^\gamma},1]$.
\end{proof}

\appendix

\section{Finite-volume reduction and locking thresholds}
\label{app:mz-reduction}

This appendix adapts the finite-volume graphical representation,
range--duration identity, and locking-threshold argument of Meester and
Znamenski \cite{MZ04}. We spell out the finite--infinite volume coupling
and the regenerative estimates in the common notation needed for both
update rules. Fix $\gamma\in\Gamma$ and retain the notation of
Section~\ref{sec:stationary-limits}.

\subsection{Avalanche coupling and durations}

Translation invariance makes the law of a finite-volume avalanche
independent of its origin. Lemma~\ref{lem:level-regeneration} remains
valid for both update rules and on $\Lambda_N$. Conditionally on the
information revealed up to level $t$, the terminal fitnesses on the
avalanche range are independent with law $t+\operatorname{Exp}(1)$.

\begin{lemma}[Finite--infinite volume coupling]
\label{lem:app-finite-infinite}
For every $N\geq2$ and $t\geq0$, the finite- and infinite-volume
$\gamma$-avalanches may be coupled so that
\[
        r_N^\gamma(t)=r_\infty^\gamma(t)\wedge N.
\]
Consequently,
\[
        R_N^\gamma(t)
        =\mathbb{E}[r_\infty^\gamma(t)\wedge N]
        \uparrow R_\infty^\gamma(t),
        \qquad
        P_N^\gamma(t)
        =\mathbb{P}(r_\infty^\gamma(t)\geq N)
        \downarrow P_\infty^\gamma(t).
\]
On $\{r_\infty^\gamma(t)<N\}$ the two avalanches have the same duration.
\end{lemma}

\begin{proof}
Couple the systems with the same refresh variables and project the
infinite-volume process onto $\Lambda_N$. While its range has fewer than
$N$ sites, it is an interval on which the projection is injective, so the
evolutions agree. If it stops before reaching $N$ sites, the ranges and
durations coincide. Once it reaches $N$ sites, its interval projects onto
all of $\Lambda_N$, and the finite-volume range is $N$. The limit
statements follow from monotone convergence and continuity from above.
\end{proof}



For $s\geq0$, set
$\nu_{N,s}^\gamma
:=\mathcal L(\xi_N^\gamma(s),\eta_N^\gamma(s))$, the joint law of the
range set and duration of an $s$-avalanche on $\Lambda_N$ rooted at the
origin. Independently for $x\in\Lambda_N$, take a marked Poisson point
process
\[
        \Pi_{N,x}^\gamma
        =
        \sum_{j\geq1}
        \delta_{(\tau_{N,x,j}^\gamma,
        A_{N,x,j}^\gamma,h_{N,x,j}^\gamma)}
        \quad\text{with intensity }
        ds\,\nu_{N,s}^\gamma(dA,dh)
\]
on $(0,\infty)\times2^{\Lambda_N}\times\N$. Here
$\tau_{N,x,j}^\gamma$ is the level of the mark,
$A_{N,x,j}^\gamma$ is the range set of the attached avalanche relative
to its origin, and $h_{N,x,j}^\gamma$ is its duration. 


At a mark $(x,s,A,h)$, draw horizontal arrows from $(x,s)$ to
$(x+y,s)$ for every $y\in A$, with addition modulo $N$, and allow
upward motion along the vertical lines. Starting from
$(\mathcal N_\gamma,0)$, with the seed projected onto $\Lambda_N$ and
repetitions removed, let $\xi_N^{\gamma,\mathrm{gr}}(t)$ be the set of
sites reachable at level $t$, and define
\[
        \eta_N^{\gamma,\mathrm{gr}}(t)
        :=
        1+
        \sum_{x\in\Lambda_N}
        \sum_{\tau_{N,x,j}^\gamma\leq t}
        \mathbf 1_{\{
        (x,\tau_{N,x,j}^\gamma)\text{ is reachable}\}}
        h_{N,x,j}^\gamma.
\]
The initial term counts the forced update, while every reachable mark
contributes the duration of the subavalanche attached to it. Since
$\Lambda_N$ is finite and every vertical line carries a rate-one
Poisson process, almost surely there are only finitely many marks below
each fixed level.

\begin{lemma}[Finite-volume graphical representation]
\label{lem:app-graphical}
The graphical and physical avalanche families have the same law:
\[
        \bigl(
        \xi_N^{\gamma,\mathrm{gr}}(t),
        \eta_N^{\gamma,\mathrm{gr}}(t)
        \bigr)_{t\geq0}
        \stackrel{\mathrm d}=
        \bigl(
        \xi_N^\gamma(t),
        \eta_N^\gamma(t)
        \bigr)_{t\geq0}.
\]
In particular, their range and duration are non-decreasing in the
level.
\end{lemma}

\begin{proof}
Explore the reachable marks in increasing order of level. Suppose that
the marks below the current level have been explored and that the
current reachable set is $I\subseteq\Lambda_N$. The next mark carried
by a vertical line through $I$ occurs at total rate $|I|$, its base
point is uniform on $I$, and, conditionally on its level $s$, its mark
$(A,h)$ is independent of the explored past with law
$\nu_{N,s}^\gamma$.

By level regeneration, this is exactly the conditional law of the next
subavalanche appearing when the threshold of the physical avalanche is
raised: its origin is uniform on the current range, and its translated
range and duration are those of an independent $s$-avalanche.
Induction over the finitely many reachable marks below each fixed level
identifies the two families, while monotonicity follows directly from
reachability.
\end{proof}



\begin{proposition}[Range--duration identity]
\label{prop:app-range-duration}
For every $N\geq2$ and $t\geq0$,
\begin{equation}
\label{eq:app-finite-duration}
        D_N^\gamma(t)
        =
        1+\int_0^tR_N^\gamma(s)D_N^\gamma(s)\,ds
        =
        \exp\left\{\int_0^tR_N^\gamma(s)\,ds\right\}.
\end{equation}
If $R_\infty^\gamma(t)<\infty$, then
\begin{equation}
\label{eq:app-infinite-duration}
        D_\infty^\gamma(t)
        =
        \exp\left\{\int_0^tR_\infty^\gamma(s)\,ds\right\}<\infty.
\end{equation}
Consequently,
\[
        t_{\mathrm d}^\gamma=t_{\mathrm{sus}}^\gamma,
\]
and, for every $t<t_{\mathrm{sus}}^\gamma$,
\begin{equation}
\label{eq:app-duration-convergence}
        D_N^\gamma(t)\longrightarrow D_\infty^\gamma(t).
\end{equation}
\end{proposition}

\begin{proof}
For fixed $N$ and $t<\infty$, $\eta_N^\gamma(t)$ is integrable. Indeed,
during any block of $N$ updates, the event that every refreshed value
exceeds $t$ has conditional probability at least $e^{-d_\gamma tN}$; on
this event no new value below $t$ is created and the selected minimum
removes at least one such value at each update, so the avalanche stops
within the block.

By Lemma~\ref{lem:app-graphical},
\[
        \eta_N^\gamma(t)
        =
        1+
        \sum_{x\in\Lambda_N}
        \sum_{\tau_{N,x,j}^\gamma\leq t}
        \mathbf 1_{\{x\in
        \xi_N^\gamma(\tau_{N,x,j}^\gamma-)\}}
        h_{N,x,j}^\gamma.
\]
The indicator is predictable, and conditionally on a mark at height $s$
its duration is independent of the past with mean $D_N^\gamma(s)$. The
compensation formula gives the first equality in
\eqref{eq:app-finite-duration}. Moreover,
\[
        0\leq
        R_N^\gamma(s+h)-R_N^\gamma(s)
        \leq N(1-e^{-Nh}),
\]
since the range cannot change if no active vertical line carries a mark
between levels $s$ and $s+h$. Thus $R_N^\gamma$ is continuous, and the
scalar integral equation gives the exponential formula.

Assume that $R_\infty^\gamma(t)<\infty$. For every $s\leq t$ the
$s$-avalanche has finite range almost surely; the block argument used
for finite-volume integrability then also shows that its duration is
finite almost surely. In the infinite-volume graphical construction,
decorate each mark at level $u$ with an independent copy of the full
genealogical $u$-avalanche, including its duration. Let
$\eta_\infty^{\gamma,[m]}(s)$ be the duration obtained by retaining the
initial forced update and at most $m$ nested graft generations. Then
\[
        \eta_\infty^{\gamma,[0]}(s)=1,
        \qquad
        \eta_\infty^{\gamma,[m]}(s)
        \uparrow\eta_\infty^\gamma(s).
\]

Set
$D_\infty^{\gamma,[m]}(s)
:=\mathbb{E}[\eta_\infty^{\gamma,[m]}(s)]$. At level $u$, the expected
number of active base points is $R_\infty^\gamma(u)$, while the
truncated duration attached to each mark is independent of the explored
past and has mean $D_\infty^{\gamma,[m]}(u)$. Compensation therefore
gives
\[
        D_\infty^{\gamma,[m+1]}(s)
        =
        1+\int_0^s
        R_\infty^\gamma(u)D_\infty^{\gamma,[m]}(u)\,du.
\]
Writing $H(s):=\int_0^sR_\infty^\gamma(u)\,du$, induction yields
\[
        D_\infty^{\gamma,[m]}(s)
        =
        \sum_{j=0}^m\frac{H(s)^j}{j!}.
\]
Monotone convergence now proves
\eqref{eq:app-infinite-duration}. Conversely,
$r_\infty^\gamma(t)\leq d_\gamma\eta_\infty^\gamma(t)$, so finite mean
duration implies finite mean range. Hence
$t_{\mathrm d}^\gamma=t_{\mathrm{sus}}^\gamma$.

Finally, Lemma~\ref{lem:app-finite-infinite} gives
$R_N^\gamma(s)\uparrow R_\infty^\gamma(s)$ for every $s$, and for
$s\leq t<t_{\mathrm{sus}}^\gamma$ these functions are bounded by
$R_\infty^\gamma(t)$. Dominated convergence in the exponential formulas
proves \eqref{eq:app-duration-convergence}.
\end{proof}

\subsection{Locking thresholds and the stationary law}

For $i\in\Lambda_N$, write
\[
        \mathcal U_{N,\gamma}(i)
        :=i+\mathcal N_\gamma
\]
modulo $N$, with repetitions removed, and let
\[
        F_N^\gamma(n)
        =(F_{N,x}^\gamma(n))_{x\in\Lambda_N}
\]
be the finite-volume fitness chain. Ties, which can occur only in a prescribed
deterministic initial configuration, are resolved by a fixed deterministic
rule.

\begin{proposition}[Finite-volume ergodicity]
\label{prop:app-ergodicity}
For every $N\geq2$, the $\gamma$-chain has a unique stationary law
$\pi_N^\gamma$. From every deterministic initial configuration, its law
converges to $\pi_N^\gamma$ in total variation.
\end{proposition}


\begin{proof}
If $N\leq d_\gamma$, every update refreshes the whole system, so after
one step the configuration is a collection of independent exponential
variables. The conclusion is immediate. Assume henceforth that
$N>d_\gamma$, and fix $s>0$.

From every initial configuration, the chain reaches the set
$\{F_x>s\text{ for all }x\in\Lambda_N\}$ in finite mean time. Indeed,
in each block of $N$ updates, with conditional probability at least
$e^{-d_\gamma sN}$, every refreshed value exceeds $s$. On this event
no new value below $s$ is created, while each update removes the
current minimum; hence all values below $s$ disappear within the block.

Starting from a configuration above $s$, the successive $s$-avalanches
have iid range--duration pairs. Moreover, $P_N^\gamma(s)>0$, since a
prescribed finite sequence of refresh values can force an avalanche to
visit every site. Thus a spanning $s$-avalanche occurs after a
geometric number of trials. Each $s$-avalanche has finite mean duration:
the preceding block argument, applied while the avalanche is running,
gives a geometric tail on the scale of $N$. Hence the terminal time of
the first spanning $s$-avalanche has finite mean.

At such a terminal time every site has been refreshed, and level
regeneration gives the configuration the law
\[
        \rho_{N,s}:=
        (s+\operatorname{Exp}(1))^{\otimes\Lambda_N},
\]
independently of the past. The successive terminal times of spanning
$s$-avalanches are therefore regeneration times. After the first one,
the cycles are iid and have finite mean length.

Their length distribution has span one. Indeed, for some $m\geq1$ a
spanning $s$-avalanche of duration $m$ has positive probability. It
may instead be preceded by a one-step nonspanning $s$-avalanche, which
has probability $e^{-d_\gamma s}>0$: all refreshed values then exceed
$s$, and $N>d_\gamma$ prevents spanning in one step. Successive
$s$-avalanches have the same iid range--duration law, so cycle lengths
$m$ and $m+1$ both occur with positive probability.


By the regenerative limit theorem
\cite[Theorem~VI.1.2 and Corollary~VI.1.5, pp.~170--171]{Asmussen2003},
there is a stationary regenerative law $\pi_N^\gamma$, and the
finite-mean, span-one cycle structure gives
\[
        \bigl\|
        \mathcal L(F_N^\gamma(n))-\pi_N^\gamma
        \bigr\|_{\mathrm{TV}}
        \longrightarrow0 .
\]


% The regenerative convergence theorem
% \cite[Chapter~VI]{Asmussen2003} now yields a unique stationary law
% $\pi_N^\gamma$ and, since the first regeneration time has finite mean
% uniformly in the initial configuration,
% \[
%         \bigl\|
%         \mathcal L(F_N^\gamma(n))-\pi_N^\gamma
%         \bigr\|_{\mathrm{TV}}
%         \longrightarrow0
% \]
% from every deterministic initial configuration. 
% Finally, the same
% regeneration bound applies from any initial law. If $\nu$ is stationary,
% the regenerative theorem gives
% $\mathcal L_\nu(F_N^\gamma(n))\to\pi_N^\gamma$ in total variation,
% whereas stationarity gives
% $\mathcal L_\nu(F_N^\gamma(n))=\nu$ for every $n$. Hence
% $\nu=\pi_N^\gamma$, proving uniqueness.
\end{proof}




% \begin{proof}
% If $N\leq d_\gamma$, every update refreshes the whole system, so after one
% step the configuration consists of independent exponential variables and
% the assertion is immediate. Assume $N>d_\gamma$ and fix $0<t<c$. From
% any configuration, the chain reaches a state with all fitnesses above $t$
% in finite mean time. Indeed, in any block of $N$ updates, with conditional
% probability $e^{-d_\gamma tN}$ every refreshed value exceeds $t$; on this
% event no new value below $t$ is created and every update removes the
% current minimum, so all such values disappear within the block.

% Once all fitnesses exceed $t$, successive $t$-avalanches use disjoint
% refresh variables, and their range--duration pairs are iid. A prescribed
% finite sequence of refresh values can make an avalanche visit every site,
% so $P_N^\gamma(t)>0$ and a spanning avalanche occurs after a geometric
% number of trials. At its terminal time every site has been refreshed, and
% level regeneration gives independent terminal values with law
% $t+\operatorname{Exp}(1)$. Continue the chain until all fitnesses exceed
% $c$. The lack-of-memory property then gives a configuration of independent
% $c+\operatorname{Exp}(1)$ variables, independent of the past, and the
% preceding block estimate with $t$ in place of $s$ shows that the recovery
% time has finite mean.

% At each recovery time the configuration has the fixed law of independent
% $c+\operatorname{Exp}(1)$ variables and is independent of the past; the
% successive cycles are therefore iid and integrable. Their length distribution
% has span one: a fixed positive-probability spanning-and-recovery scenario may
% be preceded by a one-step nonspanning $t$-avalanche, an event of probability
% $e^{-d_\gamma t}>0$. The discrete aperiodic renewal theorem for regenerative
% processes then gives total-variation convergence to the stationary cycle law.
% Since the same regeneration law is accessible from every initial
% configuration, the stationary law is unique.
% \end{proof}

Assume henceforth that the initial fitnesses are independent
$\operatorname{Exp}(1)$ variables. Set $Y_{N,x}^\gamma(0)=0$. If
$i_N^\gamma(n)$ is the site of minimum fitness and
\[
        M_N^\gamma(n)
        :=F_{N,i_N^\gamma(n)}^\gamma(n),
\]
define
\begin{equation}
\label{eq:app-locking-recursion}
        Y_{N,x}^\gamma(n+1)
        =
        \begin{cases}
        0,
        &x\in\mathcal U_{N,\gamma}(i_N^\gamma(n)),\\
        \max\{Y_{N,x}^\gamma(n),M_N^\gamma(n)\},
        &x\notin\mathcal U_{N,\gamma}(i_N^\gamma(n)).
        \end{cases}
\end{equation}
For $y,z\geq0$, put
\[
        H_y(z)
        :=
        \begin{cases}
        0,&z<y,\\
        1-e^{-(z-y)},&z\geq y.
        \end{cases}
\]

\begin{proposition}[Conditional exponential representation]
\label{prop:app-locking-representation}
Conditionally on $Y_N^\gamma(n)$, the fitnesses
$\{F_{N,x}^\gamma(n):x\in\Lambda_N\}$ are independent, and
$F_{N,x}^\gamma(n)$ has the exponential distribution conditioned to lie
above $Y_{N,x}^\gamma(n)$. Consequently, for distinct
$x_1,\ldots,x_k\in\Lambda_N$ and $z_1,\ldots,z_k\geq0$,
\[
        \mathbb{P}(F_{N,x_j}^\gamma(n)\leq z_j,\ 1\leq j\leq k)
        =
        \mathbb{E}\left[
        \prod_{j=1}^kH_{Y_{N,x_j}^\gamma(n)}(z_j)
        \right].
\]
\end{proposition}

\begin{proof}
The assertion holds at time $0$. Suppose it holds at time $n$, and
condition first on $Y_N^\gamma(n)$ and then on the minimum value and its
location. The non-minimal coordinates remain independent; by the
lack-of-memory property, the conditional law at a non-updated site $x$ is
exponential above
\[
        \max\{Y_{N,x}^\gamma(n),M_N^\gamma(n)\}.
\]
The updated sites receive fresh independent exponentials and have threshold
zero, exactly as in \eqref{eq:app-locking-recursion}.
\end{proof}

Write
\[
        G_N^\gamma(n,t)
        :=\mathbb{P}(Y_{N,0}^\gamma(n)\leq t);
\]
translation invariance makes the choice of the site $0$ immaterial.

\begin{proposition}[Subcritical locking thresholds]
\label{prop:app-locking-below}
For every $t<t_{\mathrm d}^\gamma$,
\[
        \lim_{N\to\infty}
        \limsup_{n\to\infty}G_N^\gamma(n,t)=0.
\]
\end{proposition}

\begin{proof}
Let
\[
        K_N^\gamma(n,t)
        :=
        \sum_{x\in\Lambda_N}
        \mathbf 1_{\{Y_{N,x}^\gamma(n)\leq t\}}.
\]
After an initial delayed cycle, let $T_j$ be the starting time of the
$j$-th $t$-avalanche and write $(r_j,\eta_j)$ for its range--duration
pair, so that $T_{j+1}-T_j=\eta_j$ and these pairs are iid. At the first
update of the avalanche the current minimum exceeds $t$, so every site
outside its eventual range has its locking threshold raised above $t$.
During the remaining updates the selected minima are at most $t$, and
only sites in the avalanche range are refreshed. Consequently,
\[
        K_N^\gamma(n,t)\leq r_j,
        \qquad T_j<n\leq T_{j+1}.
\]
The integer interval $(T_j,T_{j+1}]$ contains exactly $\eta_j$ times, so
its total reward is at most $\eta_jr_j$. For all sufficiently large
$N$, a one-step nonspanning avalanche has probability
$e^{-d_\gamma t}>0$, and the cycle length therefore has span one. The
delayed renewal-reward theorem and translation invariance give
\[
        \limsup_{n\to\infty}G_N^\gamma(n,t)
        \leq
        \frac{\mathbb{E}[\eta_N^\gamma(t)r_N^\gamma(t)]}
        {ND_N^\gamma(t)}.
\]
For $K\geq1$, splitting according to $r_N^\gamma(t)\leq K$ yields
\begin{equation}
\label{eq:app-subcritical-split}
        \limsup_{n\to\infty}G_N^\gamma(n,t)
        \leq
        \frac KN+
        \frac{
        \mathbb{E}[\eta_N^\gamma(t)
        \mathbf 1_{\{r_N^\gamma(t)>K\}}]}
        {D_N^\gamma(t)}.
\end{equation}
Set
\[
\begin{aligned}
        e_N^\gamma(t,k)
        &:=
        \mathbb{E}[\eta_N^\gamma(t)
        \mathbf 1_{\{r_N^\gamma(t)=k\}}],\\
        e_\infty^\gamma(t,k)
        &:=
        \mathbb{E}[\eta_\infty^\gamma(t)
        \mathbf 1_{\{r_\infty^\gamma(t)=k\}}].
\end{aligned}
\]
By Lemma~\ref{lem:app-finite-infinite},
$e_N^\gamma(t,k)=e_\infty^\gamma(t,k)$ for $k<N$. Thus, for fixed
$K<N$,
\[
        \mathbb{E}[\eta_N^\gamma(t)
        \mathbf 1_{\{r_N^\gamma(t)>K\}}]
        =
        D_N^\gamma(t)
        -\sum_{k=1}^Ke_\infty^\gamma(t,k).
\]
Since $t<t_{\mathrm d}^\gamma=t_{\mathrm{sus}}^\gamma$,
Proposition~\ref{prop:app-range-duration} gives
$D_N^\gamma(t)\to D_\infty^\gamma(t)<\infty$. The last display therefore
converges to $\sum_{k>K}e_\infty^\gamma(t,k)$, which tends to zero as
$K\to\infty$. Let first $N\to\infty$ in
\eqref{eq:app-subcritical-split} and then $K\to\infty$.
\end{proof}

\begin{proposition}[Supercritical locking thresholds]
\label{prop:app-locking-above}
For every $c>t_\infty^\gamma$,
\[
        \lim_{N\to\infty}
        \liminf_{n\to\infty}G_N^\gamma(n,c)=1.
\]
\end{proposition}

\begin{proof}
Fix $t_\infty^\gamma<t<t'<c$. Start at a regeneration time at
which the fitnesses are independent with law
$c+\operatorname{Exp}(1)$. Run the chain until the end of the first
spanning $t$-avalanche and denote this duration by $A_N$; then continue
until all fitnesses again exceed $c$ and denote the recovery duration by
$B_N$. 
At the end of the recovery phase all fitnesses exceed \(c\).
By the same level-regeneration argument as in Lemma~\ref{prop:app-ergodicity},
their conditional law is
\[
        (c+\operatorname{Exp}(1))^{\otimes\Lambda_N},
\]
independently of the history before that terminal time. Thus successive
avalanche--recovery periods are iid cycles. Their length has span one by
the same one-step-delay argument as in Proposition~\ref{prop:app-ergodicity}.

% Repetition gives iid cycles, whose length has span one by the same
% one-step-delay argument as in Proposition~\ref{prop:app-ergodicity}.



At the end of the spanning $t$-avalanche every site has been updated, so
all locking thresholds are at most $t$. During the recovery all selected
minima are below $c$, hence all locking thresholds remain at most $c$.
The renewal-reward theorem gives
\begin{equation}
\label{eq:app-recovery-fraction}
        \liminf_{n\to\infty}G_N^\gamma(n,c)
        \geq
        \frac{\mathbb{E}[B_N]}{\mathbb{E}[A_N]+\mathbb{E}[B_N]}.
\end{equation}
The chain runs through iid $t$-avalanches until the first spanning one, so
Wald's identity yields
\begin{equation}
\label{eq:app-avalanche-phase}
        \mathbb{E}[A_N]
        =
        \frac{D_N^\gamma(t)}{P_N^\gamma(t)}.
\end{equation}
At the start of the recovery phase the active set is all of $\Lambda_N$.
In the graphical construction, for each $x\in\Lambda_N$ let $\tau_x$ be
the first mark of $\Pi_{N,x}^\gamma$ above $t'$. Whenever
$\tau_x\leq c$, this mark is reachable by the vertical path from $(x,t)$,
and its attached duration is one of the non-negative contributions to the
recovery duration $B_N$. Conditionally on $\tau_x$, its mean is
$D_N^\gamma(\tau_x)\geq D_N^\gamma(t')$. Therefore
\[
        \mathbb{E}[B_N]
        \geq
        N(1-e^{-(c-t')})D_N^\gamma(t').
\]
Combining this estimate with
\eqref{eq:app-recovery-fraction}--\eqref{eq:app-avalanche-phase} and using
$D_N^\gamma(t')\geq D_N^\gamma(t)$ gives
\[
        1-\liminf_{n\to\infty}G_N^\gamma(n,c)
        \leq
        \frac{1}
        {N(1-e^{-(c-t')})P_N^\gamma(t)}.
\]
By Lemma~\ref{lem:app-finite-infinite},
$P_N^\gamma(t)\geq P_\infty^\gamma(t)>0$, so the right-hand side tends
to zero.
\end{proof}

\begin{proposition}[Stationary product limit]
\label{prop:app-stationary-limit}
If
\[
        t_{\mathrm d}^\gamma
        =t_{\mathrm{sus}}^\gamma
        =t_\infty^\gamma
        =:t_c^\gamma,
\]
then, for every fixed collection of distinct sites
$x_1,\ldots,x_k\in\Z$, the corresponding marginal of $\pi_N^\gamma$
converges to $\nu_{t_c^\gamma}^{\otimes k}$.
\end{proposition}

\begin{proof}
Fix $0<\delta<t_c^\gamma$. For all sufficiently large $N$, the sites
$x_1,\ldots,x_k$ are distinct modulo $N$. By translation invariance, the
union bound, and Propositions~\ref{prop:app-locking-below}
and~\ref{prop:app-locking-above},
\begin{align*}
&\lim_{N\to\infty}\limsup_{n\to\infty}
\mathbb{P}\left(
\max_{1\leq j\leq k}
|Y_{N,x_j}^\gamma(n)-t_c^\gamma|>\delta
\right)\\
&\quad\leq
k\lim_{N\to\infty}\limsup_{n\to\infty}
G_N^\gamma(n,t_c^\gamma-\delta)\\
&\qquad+
k\lim_{N\to\infty}\limsup_{n\to\infty}
\bigl(1-G_N^\gamma(n,t_c^\gamma+\delta)\bigr)=0.
\end{align*}
For $z_1,\ldots,z_k\geq0$, define
\[
        \Phi(y_1,\ldots,y_k)
        :=
        \prod_{j=1}^kH_{y_j}(z_j).
\]
This function is bounded and continuous. By
Proposition~\ref{prop:app-locking-representation},
\[
        \mathbb{P}(F_{N,x_j}^\gamma(n)\leq z_j,\ 1\leq j\leq k)
        =
        \mathbb{E}[
        \Phi(Y_{N,x_1}^\gamma(n),\ldots,Y_{N,x_k}^\gamma(n))
        ].
\]
If
\[
        \omega_\delta
        :=
        \sup_{\max_j|y_j-t_c^\gamma|\leq\delta}
        |
        \Phi(y_1,\ldots,y_k)
        -\Phi(t_c^\gamma,\ldots,t_c^\gamma)
        |,
\]
then $\omega_\delta\to0$ as $\delta\downarrow0$, and the absolute
difference between the preceding probability and
$\prod_{j=1}^kH_{t_c^\gamma}(z_j)$ is at most
\[
        \omega_\delta+
        \mathbb{P}\left(
        \max_j|Y_{N,x_j}^\gamma(n)-t_c^\gamma|>\delta
        \right).
\]
For fixed $N$, Proposition~\ref{prop:app-ergodicity} allows us to let
$n\to\infty$ and replace the law of $F_N^\gamma(n)$ by
$\pi_N^\gamma$. Letting next $N\to\infty$ and finally
$\delta\downarrow0$ proves convergence of the joint distribution
functions to
\[
        \prod_{j=1}^kH_{t_c^\gamma}(z_j),
\]
the distribution function of $\nu_{t_c^\gamma}^{\otimes k}$.
\end{proof}

\begin{thebibliography}{99}

\bibitem{Asmussen2003}
S.~Asmussen,
\emph{Applied Probability and Queues},
second ed., Stochastic Modelling and Applied Probability, vol.~51,
Springer-Verlag, New York, 2003.

\bibitem{BS93}
P.~Bak and K.~Sneppen,
\emph{Punctuated equilibrium and criticality in a simple model of evolution},
Phys. Rev. Lett. \textbf{71} (1993), 4083--4086.

\bibitem{GMN06}
A.~Gillett, R.~Meester and M.~Nuyens,
\emph{Bounds for avalanche critical values of the Bak--Sneppen model},
Markov Process. Related Fields \textbf{12} (2006), no.~4, 679--694.

\bibitem{GMW06}
A.~J.~Gillett, R.~W.~J.~Meester and P.~van der Wal,
\emph{Maximal avalanches in the Bak--Sneppen model},
J. Appl. Probab. \textbf{43} (2006), no.~3, 840--851.

\bibitem{GKW11}
M.~Grinfeld, P.~A.~Knight and A.~R.~Wade,
\emph{Bak--Sneppen-type models and rank-driven processes},
Phys. Rev. E \textbf{84} (2011), 041123.

\bibitem{GKW12}
M.~Grinfeld, P.~A.~Knight and A.~R.~Wade,
\emph{Rank-driven Markov processes},
J. Stat. Phys. \textbf{146} (2012), no.~2, 378--407.

\bibitem{Maslov96}
S.~Maslov,
\emph{Infinite series of exact equations in the Bak--Sneppen model of
biological evolution},
Phys. Rev. Lett. \textbf{77} (1996), 1182--1185.

\bibitem{MDRMZ98}
S.~Maslov, P.~De Los Rios, M.~Marsili and Y.-C.~Zhang,
\emph{Critical exponents of the anisotropic Bak--Sneppen model},
Phys. Rev. E \textbf{58} (1998), 7141--7145.

\bibitem{MZ03}
R.~Meester and D.~Znamenski,
\emph{Limit behavior of the Bak--Sneppen evolution model},
Ann. Probab. \textbf{31} (2003), 1986--2002.

\bibitem{MZ04}
R.~Meester and D.~Znamenski,
\emph{Critical thresholds and the limit distribution in the Bak--Sneppen
model},
Comm. Math. Phys. \textbf{246} (2004), 63--86.

\bibitem{PMB96}
M.~Paczuski, S.~Maslov and P.~Bak,
\emph{Avalanche dynamics in evolution, growth, and depinning models},
Phys. Rev. E \textbf{53} (1996), 414--443.

\bibitem{Volkov22}
S.~Volkov,
\emph{Rigorous upper bound for the discrete Bak--Sneppen model},
J. Stat. Phys. \textbf{186} (2022), article~8.

\end{thebibliography}

\end{document}
